% Options for packages loaded elsewhere
\PassOptionsToPackage{unicode}{hyperref}
\PassOptionsToPackage{hyphens}{url}
\documentclass[
]{article}
\usepackage{xcolor}
\usepackage[margin=0.85in]{geometry}
\usepackage{amsmath,amssymb}
\setcounter{secnumdepth}{5}
\usepackage{iftex}
\ifPDFTeX
  \usepackage[T1]{fontenc}
  \usepackage[utf8]{inputenc}
  \usepackage{textcomp} % provide euro and other symbols
\else % if luatex or xetex
  \usepackage{unicode-math} % this also loads fontspec
  \defaultfontfeatures{Scale=MatchLowercase}
  \defaultfontfeatures[\rmfamily]{Ligatures=TeX,Scale=1}
\fi
\usepackage{lmodern}
\ifPDFTeX\else
  % xetex/luatex font selection
\fi
% Use upquote if available, for straight quotes in verbatim environments
\IfFileExists{upquote.sty}{\usepackage{upquote}}{}
\IfFileExists{microtype.sty}{% use microtype if available
  \usepackage[]{microtype}
  \UseMicrotypeSet[protrusion]{basicmath} % disable protrusion for tt fonts
}{}
\makeatletter
\@ifundefined{KOMAClassName}{% if non-KOMA class
  \IfFileExists{parskip.sty}{%
    \usepackage{parskip}
  }{% else
    \setlength{\parindent}{0pt}
    \setlength{\parskip}{6pt plus 2pt minus 1pt}}
}{% if KOMA class
  \KOMAoptions{parskip=half}}
\makeatother
\usepackage{color}
\usepackage{fancyvrb}
\newcommand{\VerbBar}{|}
\newcommand{\VERB}{\Verb[commandchars=\\\{\}]}
\DefineVerbatimEnvironment{Highlighting}{Verbatim}{commandchars=\\\{\}}
% Add ',fontsize=\small' for more characters per line
\newenvironment{Shaded}{}{}
\newcommand{\AlertTok}[1]{\textcolor[rgb]{1.00,0.00,0.00}{\textbf{#1}}}
\newcommand{\AnnotationTok}[1]{\textcolor[rgb]{0.38,0.63,0.69}{\textbf{\textit{#1}}}}
\newcommand{\AttributeTok}[1]{\textcolor[rgb]{0.49,0.56,0.16}{#1}}
\newcommand{\BaseNTok}[1]{\textcolor[rgb]{0.25,0.63,0.44}{#1}}
\newcommand{\BuiltInTok}[1]{\textcolor[rgb]{0.00,0.50,0.00}{#1}}
\newcommand{\CharTok}[1]{\textcolor[rgb]{0.25,0.44,0.63}{#1}}
\newcommand{\CommentTok}[1]{\textcolor[rgb]{0.38,0.63,0.69}{\textit{#1}}}
\newcommand{\CommentVarTok}[1]{\textcolor[rgb]{0.38,0.63,0.69}{\textbf{\textit{#1}}}}
\newcommand{\ConstantTok}[1]{\textcolor[rgb]{0.53,0.00,0.00}{#1}}
\newcommand{\ControlFlowTok}[1]{\textcolor[rgb]{0.00,0.44,0.13}{\textbf{#1}}}
\newcommand{\DataTypeTok}[1]{\textcolor[rgb]{0.56,0.13,0.00}{#1}}
\newcommand{\DecValTok}[1]{\textcolor[rgb]{0.25,0.63,0.44}{#1}}
\newcommand{\DocumentationTok}[1]{\textcolor[rgb]{0.73,0.13,0.13}{\textit{#1}}}
\newcommand{\ErrorTok}[1]{\textcolor[rgb]{1.00,0.00,0.00}{\textbf{#1}}}
\newcommand{\ExtensionTok}[1]{#1}
\newcommand{\FloatTok}[1]{\textcolor[rgb]{0.25,0.63,0.44}{#1}}
\newcommand{\FunctionTok}[1]{\textcolor[rgb]{0.02,0.16,0.49}{#1}}
\newcommand{\ImportTok}[1]{\textcolor[rgb]{0.00,0.50,0.00}{\textbf{#1}}}
\newcommand{\InformationTok}[1]{\textcolor[rgb]{0.38,0.63,0.69}{\textbf{\textit{#1}}}}
\newcommand{\KeywordTok}[1]{\textcolor[rgb]{0.00,0.44,0.13}{\textbf{#1}}}
\newcommand{\NormalTok}[1]{#1}
\newcommand{\OperatorTok}[1]{\textcolor[rgb]{0.40,0.40,0.40}{#1}}
\newcommand{\OtherTok}[1]{\textcolor[rgb]{0.00,0.44,0.13}{#1}}
\newcommand{\PreprocessorTok}[1]{\textcolor[rgb]{0.74,0.48,0.00}{#1}}
\newcommand{\RegionMarkerTok}[1]{#1}
\newcommand{\SpecialCharTok}[1]{\textcolor[rgb]{0.25,0.44,0.63}{#1}}
\newcommand{\SpecialStringTok}[1]{\textcolor[rgb]{0.73,0.40,0.53}{#1}}
\newcommand{\StringTok}[1]{\textcolor[rgb]{0.25,0.44,0.63}{#1}}
\newcommand{\VariableTok}[1]{\textcolor[rgb]{0.10,0.09,0.49}{#1}}
\newcommand{\VerbatimStringTok}[1]{\textcolor[rgb]{0.25,0.44,0.63}{#1}}
\newcommand{\WarningTok}[1]{\textcolor[rgb]{0.38,0.63,0.69}{\textbf{\textit{#1}}}}
\usepackage{longtable,booktabs,array}
\usepackage{caption}
\captionsetup[table]{skip=6pt}
\newcounter{none} % for unnumbered tables
\usepackage{calc} % for calculating minipage widths
% Correct order of tables after \paragraph or \subparagraph
\usepackage{etoolbox}
\makeatletter
\patchcmd\longtable{\par}{\if@noskipsec\mbox{}\fi\par}{}{}
\makeatother
% Allow footnotes in longtable head/foot
\IfFileExists{footnotehyper.sty}{\usepackage{footnotehyper}}{\usepackage{footnote}}
\makesavenoteenv{longtable}
\setlength{\emergencystretch}{3em} % prevent overfull lines
\providecommand{\tightlist}{%
  \setlength{\itemsep}{0pt}\setlength{\parskip}{0pt}}
\usepackage[round,comma,sort&compress]{natbib}
\bibliographystyle{plainnat}
% ============================================================================
% REQUIRED PACKAGES - Essential for document rendering
% ============================================================================

% Mathematical typesetting (required for equations and symbols)
\usepackage{amsmath,amssymb}          % Mathematical symbols and environments
\usepackage{amsfonts}                 % Additional math fonts
\usepackage{amsthm}                   % Theorem environments

% Graphics and page layout (required for figures and formatting)
\usepackage{graphicx}                 % Include graphics (REQUIRED for figures)
\usepackage{geometry}                 % Page margins (set via pandoc --variable)
\usepackage{float}                    % Better float placement

% Tables (required for table formatting)
\usepackage{booktabs}                 % Professional tables
\usepackage{longtable}                % Long tables spanning pages
\usepackage{array}                    % Advanced table formatting

% PDF features (required for cross-references and metadata)
\usepackage{url}                      % URL formatting
\usepackage{hyperref}                 % Hyperlinks and cross-references
% natbib is loaded once by the Pandoc template; options are applied in _render_pdf_override._postprocess_tex

% ============================================================================
% PACKAGES - Improve formatting and functionality
% ============================================================================

% Table enhancements (optional but recommended)
\usepackage{multirow}                 % Multi-row table cells
\usepackage{caption}                  % caption formatting
\usepackage{subcaption}               % Sub-figures and sub-tables

% Math enhancements (optional but recommended)
\usepackage{bm}                       % Bold math symbols

% Reference enhancements (optional but recommended)
\usepackage{cleveref}                 % Intelligent cross-referencing
% Render DOIs as inline hyperlinks (not footnotes).  The doi.sty package
% is intentionally omitted because it renders DOIs as footnotes.  This
% override typesets them as clickable links within the bibliography entry.
\providecommand{\doi}[1]{\href{https://doi.org/#1}{\nolinkurl{doi:#1}}}

% Configure figure numbering and captions
\renewcommand{\figurename}{Figure}
\captionsetup{
    justification=centering,
    font=small,
    labelfont=bf,
    labelsep=period
}

% Configure table numbering and captions
\renewcommand{\tablename}{Table}
\captionsetup[table]{
    justification=centering,
    font=small,
    labelfont=bf,
    labelsep=period
}

% Configure section numbering
\setcounter{secnumdepth}{3}
\renewcommand{\thesection}{\arabic{section}}
\renewcommand{\thesubsection}{\arabic{section}.\arabic{subsection}}
\renewcommand{\thesubsubsection}{\arabic{section}.\arabic{subsection}.\arabic{subsubsection}}

% Configure equation numbering
\numberwithin{equation}{section}

% Configure hyperref for proper linking
\hypersetup{
    colorlinks=true,
    linkcolor=red,
    citecolor=red,
    urlcolor=red,
    filecolor=red,
    anchorcolor=red,
    pdfborder={0 0 0},
    bookmarks=true,
    bookmarksnumbered=true,
    bookmarkstype=toc,
    pdftitle={Ento-Linguistics: Language, Ambiguity, and Scientific Communication in Entomology},
    pdfsubject={Scientific Terminology Analysis in Entomology},
    pdfproducer={XeLaTeX}
    % pdfauthor and pdfkeywords are set dynamically in _title_page.tex
    % from config.yaml to avoid metadata drift
}

% ============================================================================
% PACKAGE CONFIGURATION
% ============================================================================

% Configure cleveref for intelligent cross-references
\crefname{section}{Section}{Sections}
\crefname{subsection}{Subsection}{Subsections}
\crefname{subsubsection}{Subsubsection}{Subsubsections}
\crefname{equation}{Equation}{Equations}
\crefname{figure}{Figure}{Figures}
\crefname{table}{Table}{Tables}
\crefname{appendix}{Appendix}{Appendices}

% Configure fonts for Unicode support with fallbacks
\usepackage{newunicodechar}
%\newunicodechar definitions for Unicode characters used in the manuscript
\newunicodechar{⁴}{\textsuperscript{4}}
\newunicodechar{₄}{\textsubscript{4}}
\newunicodechar{²}{\textsuperscript{2}}
\newunicodechar{₀}{\textsubscript{0}}
\newunicodechar{₁}{\textsubscript{1}}
\newunicodechar{₂}{\textsubscript{2}}
\newunicodechar{₃}{\textsubscript{3}}
% Arrow characters used in abstract, conclusion, and supplemental methods
\newunicodechar{→}{\ensuremath{\rightarrow}}
\newunicodechar{←}{\ensuremath{\leftarrow}}
\newunicodechar{↔}{\ensuremath{\leftrightarrow}}
% Additional symbols
\newunicodechar{≥}{\ensuremath{\geq}}
\newunicodechar{≤}{\ensuremath{\leq}}
\newunicodechar{≠}{\ensuremath{\neq}}
\newunicodechar{…}{\ldots}

% ============================================================================
% FONTS AND TYPOGRAPHY
% ============================================================================

% Use standard fonts for better compatibility
% NOTE: T1 fontenc is intentionally omitted — it is incompatible with xelatex
% (which uses Unicode natively). Including it generates font warnings that
% can cascade into compilation failures.
\usepackage{lmodern}

% ============================================================================
% CODE BLOCK STYLING
% ============================================================================

% code block styling for better contrast and readability
\usepackage{fancyvrb}
\usepackage{fvextra}   % Extends fancyvrb: adds bgcolor, breaklines, etc.
\usepackage{xcolor}
\usepackage{listings}

% Define custom colors for code blocks
\definecolor{codebg}{RGB}{248, 248, 248}      % Very light gray background
\definecolor{codeborder}{RGB}{200, 200, 200}  % Medium gray border
\definecolor{codefg}{RGB}{34, 34, 34}         % Dark gray text
\definecolor{commentcolor}{RGB}{102, 102, 102} % Comment color
\definecolor{keywordcolor}{RGB}{0, 0, 0}       % Keyword color
\definecolor{stringcolor}{RGB}{0, 102, 0}      % String color

% Configure Verbatim environment for inline code
\DefineVerbatimEnvironment{Verbatim}{Verbatim}{%
    fontsize=\small,
    frame=single,
    framerule=0.5pt,
    framesep=3pt,
    rulecolor=\color{codeborder},
    bgcolor=codebg
}

% Configure code block styling
\DefineVerbatimEnvironment{Highlighting}{Verbatim}{%
    fontsize=\footnotesize,
    frame=single,
    framerule=0.5pt,
    framesep=5pt,
    rulecolor=\color{codeborder},
    bgcolor=codebg
}

% Style inline code with \texttt
\renewcommand{\texttt}[1]{%
    \colorbox{codebg}{\color{codefg}\ttfamily #1}%
}

% Configure listings package for code blocks
\lstset{
    backgroundcolor=\color{codebg},
    basicstyle=\footnotesize\ttfamily\color{codefg},
    breakatwhitespace=false,
    breaklines=true,
    captionpos=b,
    commentstyle=\color{commentcolor},
    deletekeywords={...},
    escapeinside={\%*}{*)},
    extendedchars=true,
    frame=single,
    framerule=0.5pt,
    framesep=5pt,
    keepspaces=true,
    keywordstyle=\color{keywordcolor}\bfseries,
    language=Python,
    morekeywords={*,...},
    numbers=left,
    numbersep=5pt,
    numberstyle=\tiny\color{codefg},
    rulecolor=\color{codeborder},
    showspaces=false,
    showstringspaces=false,
    showtabs=false,
    stepnumber=1,
    stringstyle=\color{stringcolor},
    tabsize=4,
    title=\lstname
}

% Override any Pandoc default lstset configurations
\AtBeginDocument{
    \lstset{
        backgroundcolor=\color{codebg},
        basicstyle=\footnotesize\ttfamily\color{codefg},
        frame=single,
        framerule=0.5pt,
        framesep=5pt,
        rulecolor=\color{codeborder},
        numbers=left,
        numbersep=5pt,
        numberstyle=\tiny\color{codefg}
    }
}

% Configure bibliography
% Note: Using plainnat with natbib package for proper citation processing
% The bibliography style and commands (\bibliographystyle and \bibliography) are in 99_references.md

% Simple page break support for document structure
% Note: Page breaks are handled in the markdown generation, not here

% ============================================================================
% DOCUMENT FORMATTING
% ============================================================================

% Ensure proper spacing and formatting
\frenchspacing  % Single space after periods
\linespread{1.2}  % Slightly increased line spacing for readability

% ============================================================================
% TITLE PAGE OVERRIDE
% ============================================================================
% The title page, author block, DOI, and PDF metadata are generated at
% build time by _render_pdf_override.py from config.yaml, written to
% _title_page.tex, and \input'd here.  This avoids hardcoding metadata
% in two places.  The generated file redefines \maketitle using \parbox
% blocks instead of the default tabular (which breaks under TeXLive 2026).
\IfFileExists{_title_page.tex}{% Auto-generated from config.yaml — do not edit by hand.
% Uses a custom \maketitle to avoid TeXLive 2026 tabular/\and breakage.
\makeatletter
\renewcommand{\maketitle}{%
  \begin{center}
    {\LARGE \@title \par}
    \vskip 1.5em
    \parbox[t]{0.45\textwidth}{\centering
      \textbf{Daniel Ari Friedman}\\
      {\small Active Inference Institute, APOIDEAS}\\
      {\footnotesize {\ttfamily daniel@activeinference.institute} \enspace $\cdot$ \enspace ORCID: 0000-0001-6232-9096}
    }
    \hfill
    \parbox[t]{0.45\textwidth}{\centering
      \textbf{Tucker Cahill Chambers}\\
      {\small APOIDEAS}\\
      {\footnotesize ORCID: 0009-0008-3793-7872}
    }
    \vskip 1em
    April 15, 2026\\[4pt]
    {\small DOI: \href{https://doi.org/10.5281/zenodo.19574118}{10.5281/zenodo.19574118}}
  \end{center}
  \par\vskip 1.5em
}
\makeatother

% PDF metadata (synced from config.yaml keywords and authors)
\hypersetup{%
  pdfauthor={Daniel Ari Friedman and Tucker Cahill Chambers},
  pdfkeywords={entomology, scientific terminology, language analysis, discourse analysis, ant biology, eusociality, conceptual frameworks, active inference, terminology networks, Markov blankets, stigmergy, computational linguistics},
}
}{}
\usepackage{bookmark}
\IfFileExists{xurl.sty}{\usepackage{xurl}}{} % add URL line breaks if available
\urlstyle{same}
% fallback for those not using the hyperref driver hyperxmp:
\makeatletter
\@ifundefined{xmpquote}{\newcommand{\xmpquote}[1]{#1}}{}
\makeatother
\hypersetup{
  pdftitle={Ento-Linguistics: Language{,} Ambiguity{,} and Scientific Communication in Entomology},
  pdfauthor={Daniel Ari Friedman; Tucker Cahill Chambers},
  colorlinks=true,linkcolor=red,urlcolor=red,citecolor=red,anchorcolor=red,filecolor=red,
  pdfcreator={LaTeX via pandoc}}

\title{Ento-Linguistics: Language, Ambiguity, and Scientific
Communication in Entomology}
\usepackage{etoolbox}
\makeatletter
\providecommand{\subtitle}[1]{% add subtitle to \maketitle
  \apptocmd{\@title}{\par {\large #1 \par}}{}{}
}
\makeatother
\subtitle{How Terminology Networks Shape Understanding of Insect Biology
(And Vice-Versa)}
\author{Daniel Ari Friedman \and Tucker Cahill Chambers}
\date{2026-09-15}

\begin{document}
\maketitle

{
\setcounter{tocdepth}{2}
\tableofcontents
}
\section{Abstract}\label{sec:abstract}

Scientific language does not merely describe biological phenomena; it
actively constitutes the generative models through which researchers
parse complex systems. This paper makes three core contributions to
understanding---and correcting---the epistemic consequences of this
constitutive role. First, we introduce a six-domain Ento-Linguistic
framework that decomposes the terminological landscape of insect
research into analytically tractable themes, isolating domains where
anthropomorphic language most severely distorts causal modeling. Second,
we develop an open-source computational pipeline that integrates
automated term extraction, co-occurrence network construction, and
information-theoretic ambiguity scoring with principles from Active
Inference and Complex Systems Theory. Third, we propose and validate
four evidence-based meta-standards---Clarity, Appropriateness,
Consistency, and Evolvability (CACE)---as a formalized protocol for
lexical engineering. Analysis of a corpus encompassing 7608
entomological publications (999905 tokens; 47227 unique token types;
Type--Token Ratio 0.0472) extracts 11616 candidate terms (with 1239
assigned to specific semantic domains across 6 conceptual clusters
linked by 15 weighted relationships). The resulting terminology networks
display strong modularity alongside systematic cross-domain
bridging---most prominently in the Power and Labor domain, where 139
bridging terms generate extensive semantic bleed-over into adjacent
domains. Terms such as ``queen'' (4053 occurrences), ``worker'' (6396),
and ``caste'' (762) implicitly impose hierarchical control topologies
onto biological structures that are fundamentally stigmergic and
decentralized. Across all 1239 domain-assigned terms, 12.8\% exhibit
context-dependent semantic drift, demonstrating how conceptual
constructs like ``individuality'' span multiple biological scales and
consequently blur the formal systemic boundaries (Markov Blankets)
required for mathematically rigorous modeling. The accompanying fully
reproducible computational pipeline provides the quantitative analytical
tools necessary for a more self-aware and epistemically rigorous
scientific practice. All code and data are available at
https://github.com/docxology/ento\_linguistics.

\newpage

\section{Introduction}\label{sec:introduction}

\subsection{Linguistic Priors and Generative
Models}\label{linguistic-priors-and-generative-models}

Scientific inquiry is a process of \textbf{active inference}, where
researchers refine generative models to minimize surprise about
biological observations \cite{friston2010free}. Language acts as the
\textbf{hyper-prior} for these models: it constrains the hypothesis
space before data collection begins. When entomologists employ terms
like ``queen'' or ``caste,'' they are not merely labeling phenomena;
they are importing a high-precision prior from human social systems into
their model of insect biology. If this prior is structurally misaligned
with the target system---for instance, assuming top-down control in a
stigmergic network---the resulting model will necessarily suffer from
high variational free energy, manifesting as persistent anomalies and
theoretical epicycles \cite{kuhn1996, clark2013whatever}.

The \textbf{scientific community itself can be modeled as a multi-scale
Active Inference agent} whose collective task is to minimize long-term
surprise about the entomological world it observes. Its generative model
is the shared ontology of the field---the lexicon and conceptual
structures encoded in the literature. When this ontology is precise and
plastic, the community efficiently updates its priors in response to new
evidence (e.g., genomic data revealing that caste determination is a
labile epigenetic process rather than a fixed fate). When the ontology
is rigid or laden with hidden anthropomorphic priors, the agent suffers
from \textbf{prior dogmatism}: a failure of belief updating where
high-precision, fixed priors overwhelm contradictory sensory evidence.
In this state, anomalies are explained away rather than used to update
the model. Terminology reform is therefore a \textbf{model selection}
process: optimizing the community's generative model to restore its
capacity for free-energy minimization.

This optimization requires specific criteria. We propose
\textbf{Evolvability}---defined here as \textbf{scale-invariance}---as a
critical metric for scientific terms. An evolvable term maintains its
validity across biological scales (gene, organism, superorganism)
without fracturing. ``Queen,'' by contrast, is scale-brittle: it
functions as a metaphor at the colony level but dissolves into
incoherence when applied to the underlying genetic or molecular
mechanisms of reproductive differentiation.

The consequences of this misalignment are not merely philosophical. They
propagate through every stage of the research cycle---from hypothesis
formulation, through variable selection, to the causal explanations
offered for observed phenomena. The following section formalizes this
propagation as a problem of model integrity.

\subsection{Motivation: Minimizing Model
Misspecification}\label{motivation-minimizing-model-misspecification}

The drive for terminological clarity is not a stylistic preference but a
requirement for model integrity. As \citet{keller1991language} argued,
the language of science constitutes the cognitive scaffolding of
research. In the framework of Active Inference, an undefined or
metaphor-laden term introduces \textbf{irreducible uncertainty}
(entropy) into the scientific communication channel, degrading the
precision of the community's collective generative model.

The present moment demands this formalization. Recent cognitive science
emphasizes the distinction between deliberate and conventional metaphor
use, demonstrating that metaphor in scientific discourse often operates
as a conscious communicative strategy rather than an automatic
conceptual mapping \cite{steen2017deliberate}. Rather than perpetuating
inherited assumptions in our linguistic ontology, researchers must
critically assess whether their terminological priors minimize or
maximize the complexity of their biological models.

A paradigmatic example is the ``slave-making'' debate.
\citet{herbers2006} showed that the term ``slave'' naturalizes a human
institution while obscuring the biological mechanism of \textbf{social
parasitism}. In formal terms, the ``slave'' metaphor implies a conscious
coercion policy, whereas the replacement term ``dulosis'' correctly
identifies the phenomenon as a breakdown in nestmate recognition
signals---a failure of the Markov Blanket's security filter. Reform here
is not merely ethical; it restores the causal fidelity of the scientific
model by replacing a high-entropy metaphor with a mechanistically
precise descriptor.

\subsection{The Challenge of Terminological
Reform}\label{the-challenge-of-terminological-reform}

A common objection to terminological reform is that changing vocabulary
creates disconnection from existing literature. If entomologists abandon
terms like ``caste'' or ``slave,'' how would researchers locate papers
on task performance or social parasitism?

This objection inadvertently strengthens the case for reform. Retaining
problematic terminology for convenience perpetuates and compounds the
conceptual distortions it encodes \cite{herbers2006}. The appropriate
response is systematic development of clearer vocabulary alongside the
indexing infrastructure needed for literature
continuity---cross-referencing deprecated terms, establishing synonym
mappings, and leveraging modern search capabilities that already make
vocabulary-independent retrieval routine. Growing professional consensus
around inclusive language in myrmecology and the Entomological Society
of America's Better Common Names Project \cite{betternamesproject2024}
demonstrate that the field increasingly recognizes both the necessity
and the feasibility of reform.

\subsection{Ento-Linguistic Domains: A Framework for
Analysis}\label{ento-linguistic-domains-a-framework-for-analysis}

We organize our analysis around six domains where entomological language
creates ambiguity or imports unjustified assumptions. Each domain
isolates a distinct category of terminological friction between human
conceptual frameworks and ant biology.

\textbf{Unit of Individuality.} The definition of a biological
individual is formally equivalent to the specification of a
\textbf{Markov Blanket}---the statistical boundary separating internal
states from external states \cite{friston2013life}. Terms like
``colony,'' ``superorganism,'' and ``individual'' confuse these
boundaries, creating models where the relevant unit of agency is
undefined. Critically, the term ``colony'\,' also carries a fraught
ideological history: as \citet{vis2026colony} demonstrates, its
too-casual adoption across entomological literature imports
settler-colonial assumptions about social arrangements into descriptions
of insect life, compounding the epistemic problem of misspecified Markov
Blanket boundaries with a broader political--historical distortion.

\textbf{Behavior \& Identity.} Task performance in ants is a fluid
process of \textbf{policy selection} based on local cues
\cite{gordon2010}. However, terminology transforms these transient
policies into categorical identities (``forager,'' ``nurse''). This
effectively hard-codes a fixed-role prior into the model, obscuring the
plasticity and Bayesian updating that actually drives task allocation.

\textbf{Power \& Labor.} Terms like ``queen,'' ``worker,'' and ``caste''
impose a hierarchical control architecture on a system that is
fundamentally \textbf{stigmergic}. This introduces a causal error: it
attributes colony-level regulation to centralized agency (the queen)
rather than distributed feedback loops, fundamentally misrepresenting
the system's control theory.

\textbf{Sex \& Reproduction.} Terms like ``sex determination'' and ``sex
differentiation'' carry implicit assumptions about binary systems that
may not map onto ant reproductive biology, where haplodiploidy creates
fundamentally different patterns \cite{chandra2021epigenetics}.

\textbf{Kin \& Relatedness.} Human kinship terminology, grounded in
bilateral relatedness, creates systematic friction when applied to ant
societies structured by haplodiploidy. In haplodiploid species, full
sisters share an average relatedness coefficient of
\(r = 0.75\)---higher than the mother--daughter coefficient of
\(r = 0.5\)---a fundamental asymmetry absent from human kinship models.
Terms such as ``sister,'' ``mother,'' and ``family'' obscure this
asymmetry and its profound consequences for kin selection theory
\cite{chandra2021epigenetics}.

\textbf{Economics.} Economic metaphors---markets, trade, investment,
cost-benefit---shape analysis of ant foraging, resource distribution,
and colony-level resource management. This domain investigates how
transactional frameworks constrain biological interpretation by
conflating proximate energetic expenditure with ultimate fitness costs,
importing assumptions of rational optimisation from microeconomics into
systems that operate through evolved heuristics rather than deliberative
calculation. In Active Inference terms, economic metaphors impose a
\textbf{utility-maximising} generative model on systems that instead
minimise variational free energy through local policy selection---a
distinction with profound consequences for how foraging efficiency,
brood investment, and inter-colony resource flows are modelled and
interpreted.

\subsection{Research Approach}\label{research-approach}

This work employs a mixed-methodology framework combining computational
text analysis with theoretical discourse examination. The computational
component processes a \textbf{corpus of 7608 entomological publications}
(999905 tokens; 47227 unique token types; 11616 extracted candidate
terms, 1239 domain-assigned) using automated term extraction,
co-occurrence network construction, and information-theoretic ambiguity
scoring. The theoretical component, informed by
\citeauthor{foucault1972archaeology}'s archaeological method
\citeyearpar{foucault1972archaeology}, conceptual metaphor theory
\cite{lakoff1980metaphors}, and \citeauthor{gordon2023ecology}'s
\citeyearpar{gordon2023ecology} ecological framework for collective
behavior, examines how the statistical patterns reflect deeper
conceptual structures. Longitudinal case studies of ``caste'' and
``superorganism'' vocabularies (Section \ref{sec:experimental_results})
track terminological evolution alongside empirical discoveries over five
decades, providing diachronic evidence for the framework's claims. All
data and analysis code are reproducible and available for validation.

\subsection{Terminology Network
Visualization}\label{terminology-network-visualization}

To illustrate the framework's output, Figure \ref{fig:concept_map} shows
how terms cluster around the six Ento-Linguistic domains and form
cross-domain networks of meaning; detailed quantitative analysis follows
in Section \ref{sec:experimental_results}.

\begin{figure}[h]
\centering
\includegraphics[width=0.9\textwidth]{/Volumes/external_drive/Git/projects/ongoing/docxology/ento_linguistics/output/figures/concept_map.png}
\caption{Conceptual map of Ento-Linguistic domains showing relationships between terminology networks. Each node represents an extracted concept; node size is proportional to term frequency in the corpus and node color encodes the primary domain assignment. Edges connect co-occurring concepts, with thickness reflecting co-occurrence strength. The six domains form interconnected clusters; central hub terms such as ``colony,'' ``caste,'' and ``individual'' bridge multiple domains, demonstrating how specific terminological choices propagate across the scientific discourse of entomology.}
\label{fig:concept_map}
\end{figure}

\newpage

\section{Methods}\label{sec:methodology}

Our methodology combines two sequential phases: systematic corpus
construction through multi-source literature mining, followed by a
multi-layer computational and discourse-analytic pipeline applied to the
assembled corpus. This section describes each phase at the level of
detail required for replication; full implementation specifications are
provided in Supplemental Methods \ref{sec:supplemental_methods}, and
extended theoretical derivations in Supplemental Analysis
\ref{sec:supplemental_analysis}.

\begin{center}\rule{0.5\linewidth}{0.5pt}\end{center}

\subsection{Data Acquisition}\label{sec:data_acquisition}

\subsubsection{Search Strategy and Source
Selection}\label{search-strategy-and-source-selection}

The primary corpus was assembled by querying two open-access databases
--- \textbf{PubMed} (NCBI) and \textbf{arXiv} --- using the
\texttt{PubMedMiner} and \texttt{ArXivMiner} classes implemented in
\texttt{src/data/literature\_mining.py}. Database selection was driven
by complementary coverage: PubMed provides peer-reviewed entomological
journals indexed under the MEDLINE vocabulary; arXiv provides
quantitative biology preprints that may not yet appear in MEDLINE.

\textbf{PubMed query} (\texttt{create\_entomology\_query()}):

\begin{Shaded}
\begin{Highlighting}[]
\NormalTok{("ants" OR "Formicidae" OR "Hymenoptera" OR "eusocial" OR "eusociality"}
\NormalTok{OR "social insects" OR "colony" OR "nest" OR "foraging"}
\NormalTok{OR "division of labor") AND (English[Language])}
\end{Highlighting}
\end{Shaded}

The query was submitted to the NCBI Entrez Utilities API
(\texttt{esearch.fcgi}, \texttt{retmode=json}). Results were retrieved
in batches of 20 PMIDs via \texttt{eSummary} (metadata: title, authors,
journal, year, DOI) and \texttt{eFetch} (\texttt{rettype=abstract},
\texttt{retmode=xml}) for abstract text, with a 500 ms inter-batch delay
to comply with NCBI rate limits. \texttt{PubMedMiner} caches search
results (\texttt{enable\_cache=True}) to prevent redundant API calls
during pipeline re-runs.

\textbf{arXiv query} (\texttt{ArXivMiner.search()}):

\begin{Shaded}
\begin{Highlighting}[]
\NormalTok{cat:q{-}bio.PE OR cat:q{-}bio.QM}
\end{Highlighting}
\end{Shaded}

Results retrieved via the \href{http://export.arxiv.org/api/query}{arXiv
API} (Atom XML, \texttt{sortBy=submittedDate},
\texttt{sortOrder=descending}). Records were post-filtered by keyword
overlap against the entomology vocabulary set \{\emph{ant, ants,
formicidae, eusocial, colony, social insect}\}; only records whose
combined title+abstract text contained at least one keyword were
retained.

\subsubsection{Corpus Composition and
Cleaning}\label{corpus-composition-and-cleaning}

The assembled \texttt{LiteratureCorpus} stores \texttt{Publication}
dataclass objects with the following fields: \texttt{title},
\texttt{authors}, \texttt{abstract}, \texttt{doi}, \texttt{pmid},
\texttt{year}, \texttt{journal}, \texttt{keywords}, \texttt{full\_text}.
After deduplication (by PMID) and quality filtering (records with
neither abstract nor full text were excluded), the final corpus
comprises:

{\def\LTcaptype{none} % do not increment counter
\begin{longtable}[]{@{}ll@{}}
\toprule\noalign{}
Metric & Value \\
\midrule\noalign{}
\endhead
\bottomrule\noalign{}
\endlastfoot
Documents & \textbf{7608} \\
Total processed tokens & \textbf{999905} \\
Unique token types & \textbf{47227} \\
Candidate terms extracted & \textbf{11616} \\
Domain-assigned terms & \textbf{1239} \\
\end{longtable}
}

These statistics are computed at runtime by
\texttt{TextProcessor.get\_vocabulary\_stats()} and serialized to
\texttt{output/data/corpus\_statistics.json}; the values reported here
are read directly from that file and are therefore always current with
the last pipeline run.

Full text preprocessing --- tokenization (\texttt{sent\_tokenize},
\texttt{word\_tokenize}), scientific-term merge, stop-word removal (NLTK
English + \texttt{SCIENTIFIC\_STOP\_WORDS}), and lemmatization
(\texttt{WordNetLemmatizer}) --- is implemented in
\texttt{TextProcessor} (\texttt{src/analysis/text\_analysis.py}), with
\texttt{process\_text()} and \texttt{normalize\_text()} orchestrating
the pipeline.

\subsubsection{Domain Coverage
Verification}\label{domain-coverage-verification}

To verify that the search strategy captured all six target
Ento-Linguistic domains, term counts were computed across domains
immediately after corpus construction. The six domains and their seed
vocabularies are:

{\def\LTcaptype{none} % do not increment counter
\begin{longtable}[]{@{}
  >{\raggedright\arraybackslash}p{(\linewidth - 2\tabcolsep) * \real{0.2857}}
  >{\raggedright\arraybackslash}p{(\linewidth - 2\tabcolsep) * \real{0.7143}}@{}}
\toprule\noalign{}
\begin{minipage}[b]{\linewidth}\raggedright
Domain
\end{minipage} & \begin{minipage}[b]{\linewidth}\raggedright
Example Seed Terms
\end{minipage} \\
\midrule\noalign{}
\endhead
\bottomrule\noalign{}
\endlastfoot
Power \& Labor & queen, worker, dominance, hierarchy, division of
labor \\
Unit of Individuality & colony, superorganism, eusocial, individual,
organism \\
Sex \& Reproduction & mating, haplodiploidy, parthenogenesis, queen,
egg \\
Behavior \& Identity & caste, forager, nurse, task, polyethism \\
Kin \& Relatedness & kin selection, inclusive fitness, relatedness,
altruism \\
Economics & foraging, cost, benefit, resource allocation, trade-off \\
\end{longtable}
}

The current pipeline run extracted \textbf{11616 terms distributed
across all six domains}, of which \textbf{1239 receive specific domain
assignments}, sourced from \texttt{output/data/domain\_statistics.json}.
Domain-specific acquisition details, bridging term frequencies, and
per-domain confidence statistics are reported in Supplemental Results
\ref{sec:supplemental_results}.

\begin{center}\rule{0.5\linewidth}{0.5pt}\end{center}

\subsection{Statistical Analysis}\label{sec:statistical_analysis}

\subsubsection{Analytical Framework
Overview}\label{analytical-framework-overview}

The statistical pipeline comprises six interdependent analytical layers
applied sequentially to the assembled corpus: (1) term extraction and
classification, (2) semantic entropy estimation, (3) domain-level
statistical testing, (4) conceptual network construction and centrality
analysis, (5) rhetorical and discourse pattern scoring, and (6) CACE
meta-standard evaluation. All analyses are implemented in
\texttt{src/analysis/} and are fully deterministic
(\texttt{random\_state=42} throughout). Extended statistical derivations
and robustness tests are presented in Supplemental Analysis
\ref{sec:supplemental_analysis}.

\subsubsection{Term Extraction and
Classification}\label{term-extraction-and-classification}

\texttt{TerminologyExtractor}
(\texttt{src/analysis/term\_extraction.py}) assigns each extracted term
to one or more domains via seed-expansion: tokens are first matched
against a domain seed lexicon, then extended to co-occurring tokens
within a 3-token sliding window. Each \texttt{Term} dataclass carries
\texttt{text}, \texttt{lemma}, \texttt{domains}, \texttt{frequency},
\texttt{contexts} (deduplicated sentences), \texttt{pos\_tags},
\texttt{confidence}, and \texttt{semantic\_entropy}. N-gram extraction
(\texttt{TextProcessor.extract\_ngrams}) captures compound terms (e.g.,
\emph{division of labor}, \emph{kin selection}) that single-token
analysis would fragment. Full API documentation is in Supplemental
Methods \ref{sec:supplemental_methods}.

\subsubsection{Semantic Entropy}\label{semantic-entropy}

To quantify terminological ambiguity, we compute \textbf{Semantic
Entropy} \(H(t)\) for each term \(t\) with sufficient attestation
(\(\geq 5\) valid contexts):

\begin{equation}\label{eq:semantic_entropy}
H(t) = -\sum_{i=1}^{k} p_i \log_2 p_i \qquad \text{(bits)}
\end{equation}

where \(p_i\) is the empirical proportion of usage contexts assigned to
semantic cluster \(i\) by \(k\)-means (scikit-learn,
\texttt{random\_state=42}) over TF-IDF context vectors. The number of
clusters is set to
\(k = \max\!\bigl(2,\;\min(k_{\max},\; n-1,\; \max(2, \lfloor\!\sqrt{n}\rfloor))\bigr)\)
with \(k_{\max}=5\) and \(n = |C_t|\), ensuring \(k < n\) so that at
least some clusters contain multiple contexts and the resulting entropy
reflects genuine semantic spread rather than a degenerate uniform
assignment. Each result additionally reports the maximum attainable
entropy \(H_{\max} = \log_2 k\) and the normalized entropy
\(\hat{H}(t) = H(t)/H_{\max} \in [0,1]\), which rescales semantic spread
independently of the cluster count. Terms with \(H(t) > H^* = 2.0\)
bits---roughly corresponding to four or more equiprobable semantic
senses under uniform cluster sizes---are flagged
\texttt{is\_high\_entropy}. The threshold was calibrated against terms
of known polysemy (\emph{colony}, \emph{queen}) and specificity
(\emph{haplodiploidy}, \emph{trophallaxis}). Implementation:
\texttt{src/analysis/semantic\_entropy.py::calculate\_semantic\_entropy};
corpus-level results:
\texttt{src/analysis/semantic\_entropy.py::calculate\_corpus\_entropy}.

\subsubsection{Domain-Level Statistical
Tests}\label{domain-level-statistical-tests}

Cross-domain entropy comparisons use the following battery, all
implemented from scratch in \texttt{src/analysis/statistics.py}:

{\def\LTcaptype{none} % do not increment counter
\begin{longtable}[]{@{}
  >{\raggedright\arraybackslash}p{(\linewidth - 4\tabcolsep) * \real{0.2069}}
  >{\raggedright\arraybackslash}p{(\linewidth - 4\tabcolsep) * \real{0.3448}}
  >{\raggedright\arraybackslash}p{(\linewidth - 4\tabcolsep) * \real{0.4483}}@{}}
\toprule\noalign{}
\begin{minipage}[b]{\linewidth}\raggedright
Test
\end{minipage} & \begin{minipage}[b]{\linewidth}\raggedright
Function
\end{minipage} & \begin{minipage}[b]{\linewidth}\raggedright
Application
\end{minipage} \\
\midrule\noalign{}
\endhead
\bottomrule\noalign{}
\endlastfoot
Two-sample Welch's \(t\)-test & \texttt{t\_test} & Pairwise entropy
comparison between domains \\
One-way ANOVA & \texttt{anova\_test} & Simultaneous entropy comparison
across all 6 domains \\
95\% confidence intervals & \texttt{calculate\_confidence\_interval} &
Mean entropy uncertainty per domain \\
Pearson / Spearman correlation & \texttt{calculate\_correlation} &
Entropy--frequency relationship \\
Normal / Exponential / Uniform fit & \texttt{fit\_distribution} &
Entropy distribution characterization \\
\end{longtable}
}

The Welch--Satterthwaite degrees-of-freedom approximation is applied in
all two-sample \(t\)-tests; \(p\)-values are computed via
\texttt{scipy.stats.t.sf} and \texttt{scipy.stats.f.sf}. For the
\(\binom{6}{2} = 15\) pairwise domain comparisons, \(p\)-values are
corrected using the \textbf{Benjamini--Hochberg} false discovery rate
procedure at \(q = 0.05\). Effect sizes are reported as Cohen's \(d\)
(small: 0.2, medium: 0.5, large: 0.8 \cite{cohen1988statistical}).

A four-level \textbf{multi-scale ambiguity classification} is applied to
high-entropy terms: (1) \emph{Lexical Ambiguity} --- multiple dictionary
meanings; (2) \emph{Contextual Ambiguity} --- meaning shifts based on
research tradition (e.g., ``caste'' in classical vs.~modern entomology);
(3) \emph{Scale Ambiguity} --- meaning variation across biological
scales (\(\text{gene} \to \text{organism} \to \text{colony}\)); (4)
\emph{Temporal Ambiguity} --- historical meaning evolution traceable
across publication years. The biological-scale dimension is further
formalized through the \textbf{Markov Blanket} formalism
\cite{friston2013life}.

\subsubsection{Conceptual Network
Analysis}\label{conceptual-network-analysis}

\texttt{ConceptualMapper} (\texttt{src/analysis/conceptual\_mapping.py})
constructs a \texttt{ConceptMap} of \textbf{6 concepts}
(biological\_individuality, social\_organization, reproductive\_biology,
kinship\_systems, resource\_economics, behavioral\_ecology) linked by
\textbf{15 weighted edges}. Edge weights are overlap coefficients
(Szymkiewicz--Simpson):

\begin{equation}\label{eq:overlap_coefficient}
w_{AB} = \frac{|A \cap B|}{\min(|A|, |B|)}
\end{equation}

Composite relationship strength decomposes as:
\(\text{strength} = 0.4\,w_\text{base} + 0.3\,r_\text{term} + 0.2\,r_\text{domain} + 0.1\,\mathbb{1}_\text{hierarchical}\).

Centrality analysis uses NetworkX: degree centrality, betweenness
centrality, closeness centrality, and eigenvector centrality
(\texttt{max\_iter=1000}; \texttt{PowerIterationFailedConvergence}
fallback \(\to\) 0). Concept-level results are serialized to
\texttt{output/data/concept\_map\_summary.json}. Cross-domain bridging
terms --- appearing in \(\geq 2\) domains --- are identified with
\texttt{identify\_cross\_domain\_bridges}; the current run yields 139
bridging terms in Power \& Labor and 62 in Sex \& Reproduction.

\subsubsection{Rhetorical and Discourse
Analysis}\label{rhetorical-and-discourse-analysis}

\texttt{analyze\_rhetorical\_strategies}
(\texttt{src/analysis/rhetorical\_analysis.py}) detects four strategy
types per abstract via regex:

{\def\LTcaptype{none} % do not increment counter
\begin{longtable}[]{@{}
  >{\raggedright\arraybackslash}p{(\linewidth - 2\tabcolsep) * \real{0.4762}}
  >{\raggedright\arraybackslash}p{(\linewidth - 2\tabcolsep) * \real{0.5238}}@{}}
\toprule\noalign{}
\begin{minipage}[b]{\linewidth}\raggedright
Strategy
\end{minipage} & \begin{minipage}[b]{\linewidth}\raggedright
Detection
\end{minipage} \\
\midrule\noalign{}
\endhead
\bottomrule\noalign{}
\endlastfoot
Authority &
\texttt{\textbackslash{}(.*?20\textbackslash{}d\{2\}.*?\textbackslash{})}
--- citation parentheticals \\
Analogy &
\texttt{\textbackslash{}blike\textbackslash{}s+.*?\textbackslash{}bant}
--- ant-comparison expressions \\
Generalization &
\texttt{\textbackslash{}b(all\textbackslash{}\textbar{}every\textbackslash{}\textbar{}always\textbackslash{}\textbar{}never)\textbackslash{}s+.*?\textbackslash{}bant}
--- absolutist quantifiers \\
Anecdotal &
\texttt{\textbackslash{}b(for\ example\textbackslash{}\textbar{}such\ as\textbackslash{}\textbar{}consider\textbackslash{}\textbar{}imagine)\textbackslash{}b}
--- evidential markers \\
\end{longtable}
}

\texttt{identify\_narrative\_frameworks} classifies each abstract into
one or more of four narrative types (progress, conflict, discovery,
complexity) by keyword presence.
\texttt{score\_argumentative\_structures} decomposes argumentative
strength into claim strength, evidence quality, and reasoning coherence,
averaged to an overall score. \texttt{LinguisticFeatureExtractor}
computes anthropomorphic (\(4\) patterns), hierarchical (\(4\)
patterns), and economic (\(4\) patterns) framing densities per document.
Anthropomorphic framing indicators include five conceptual categories
--- agency, communication, social contract, cognition, and hierarchy ---
as specified in
\texttt{ConceptualMapper.detect\_anthropomorphic\_concepts()}.

\subsubsection{CACE Evaluation}\label{cace-evaluation}

Each term is scored on four bounded \([0,1]\) dimensions constituting
the \textbf{CACE} framework (\texttt{src/analysis/cace\_scoring.py}):

\begin{equation}\label{eq:cace_clarity}
\text{Clarity}(t) = \max\!\left(0,\, 1 - \frac{H(t)}{\log_2 10}\right)
\end{equation}

\begin{equation}\label{eq:cace_appropriateness}
\text{Appropriateness}(t) = 1 - \bigl[0.4\cdot\mathbb{1}_{t \in \mathcal{A}}
  + 0.1\cdot|\text{overlap}(t,\mathcal{A})|
  + 0.05\cdot\max(|\text{domains}(t)|-1,\,0)\bigr]
\end{equation}

\begin{equation}\label{eq:cace_consistency}
\text{Consistency}(t) = \bar{S}_\text{cos}(\mathbf{X}_t),\quad
  \bar{S}_\text{cos} = \frac{2}{n(n-1)}\sum_{i<j}\frac{\mathbf{x}_i \cdot \mathbf{x}_j}{\|\mathbf{x}_i\|\,\|\mathbf{x}_j\|}
\end{equation}

\begin{equation}\label{eq:cace_evolvability}
\text{Evolvability}(t) = 0.5\,\min\!\left(1,\; \frac{|\text{domains}(t)|}{3}\right) + 0.5\,\min\!\left(1,\; \frac{|S_t|}{3}\right)
\end{equation}

where \(\mathcal{A}\) is the \texttt{ANTHROPOMORPHIC\_TERMS} set
(\(\ni\) queen, king, slave, worker, soldier, nurse, \ldots);
\(\mathbb{1}_{t \in \mathcal{A}}\) indicates direct set membership (base
penalty 0.4); \(\text{overlap}(t,\mathcal{A})\) is the word-level
intersection of \(t\)'s tokens with \(\mathcal{A}\) (additional 0.1 per
overlapping word); and \(\text{domains}(t)\) counts the Ento-Linguistic
domains the term spans (additional 0.05 per domain beyond the first).
The penalty weights \((0.4, 0.1, 0.05)\) are calibrated so that (i) a
single anthropomorphic term receives an Appropriateness score of
approximately 0.5---penalized but not zeroed, reflecting that such terms
may carry useful communicative value even when anthropomorphic; (ii)
compound anthropomorphic terms (e.g., ``slave-worker'') accumulate
additional penalty proportional to their conceptual load; and (iii)
cross-domain spread contributes a smaller increment, reflecting the
empirical observation that domain bridging compounds confusion more
modestly than direct anthropomorphism. Sensitivity analysis (varying
each weight \(\pm 50\%\)) confirms that the qualitative ranking of terms
by Appropriateness is robust to the specific coefficient values: the set
of terms scoring below 0.5 remains stable across all tested
configurations. The Clarity denominator \(\log_2 10 \approx 3.32\) bits
is the \texttt{DEFAULT\_MAX\_ENTROPY} constant (corresponding to 10
equiprobable semantic senses) calibrated against terms of known polysemy
and specificity. In the Consistency equation, \(\mathbf{X}_t\) is the
TF-IDF matrix of context vectors for term \(t\), \(n = |C_t|\) is the
context count, and \(\bar{S}_\text{cos}\) is the mean pairwise cosine
similarity (high near 1 = consistent usage; near 0 = heterogeneous;
returns 0.5 when \(n < 2\)). In the Evolvability equation, the first
component scores domain breadth (terms spanning \(\geq 3\) domains
receive maximum credit), while the second scores scale breadth, where
\(S_t \subseteq \{\text{gene, cell, organism, colony, population, ecosystem}\}\)
is the subset of the six defined biological scale levels represented in
term \(t\)'s contexts; division by 3 is a calibration threshold---terms
spanning three or more scales receive maximum evolvability. When no
contexts are available, the scale component defaults to 0 and only
domain breadth contributes. Anthropomorphic terms receive a baseline
Appropriateness of \(\approx 0.40\)--\(0.60\) depending on domain
breadth---they are penalised, not zeroed. The aggregate CACE score is
the arithmetic mean of the four dimensions.
\texttt{compare\_terms\_cace} returns a ranked list for all terms.
Inter-rater reliability for qualitative CACE audits is assessed via
Cohen's \(\kappa\). Full implementation and the \texttt{CACEScore}
dataclass specification are in \texttt{src/analysis/cace\_scoring.py}.

\subsubsection{Validation and
Reproducibility}\label{validation-and-reproducibility}

Results are validated through three mechanisms: (1) \textbf{internal
consistency} --- term frequency distributions checked against semantic
entropy estimates; (2) \textbf{cross-method agreement} --- rhetorical
pattern frequencies compared with domain framing scores; (3)
\textbf{external triangulation} --- comparison against existing critical
discourse analyses of entomological literature
\cite{latour1987, longino1990}. Robustness testing (subsampling
stability, parameter sensitivity, annotation consistency) is implemented
throughout \texttt{src/analysis/}.

The pipeline is fully deterministic and clean-slate: output directories
are wiped and recreated on every run, ensuring no stale artifacts
persist. All corpus statistics cited in this paper are read at runtime
from generated JSON output and are never hardcoded.

\newpage

\section{Results: Corpus Analysis and Terminology
Networks}\label{sec:experimental_results}

\subsection{Terminology Extraction Across
Domains}\label{terminology-extraction-across-domains}

Our analysis applies the mixed-methodology framework described in
Section \ref{sec:methodology} to a corpus of entomological literature.
The dataset includes abstracts from foundational works by Hölldobler,
Wilson, and Gordon, incorporating terminology patterns characteristic of
journals including \emph{Behavioral Ecology}, \emph{Journal of Insect
Behavior}, and \emph{Insectes Sociaux}.

Domain-specific extraction from \textbf{7608 publications} (999905
tokens) identified \textbf{11616 candidate terms} total, of which
\textbf{1239 receive domain assignments} spanning all six domains, with
substantial variation in usage patterns:

\begin{table}[h]
\centering
\begin{tabular}{|l|c|c|c|}
\hline
\textbf{Domain} & \textbf{Term Count} & \textbf{Total Frequency} & \textbf{Bridging Terms} \\
\hline
Unit of Individuality & 520 & 8095 & 24 \\
Behavior \& Identity & 243 & 9239 & 91 \\
Power \& Labor & 277 & 5621 & 139 \\
Sex \& Reproduction & 227 & 4465 & 62 \\
Kin \& Relatedness & 72 & 2569 & 4 \\
Economics & 65 & 2009 & 3 \\
\hline
\end{tabular}
\caption{Domain-assigned terminology extracted from the 7608-publication corpus. Terms are assigned by seed-expansion matching against domain-specific seed vocabularies; a single term may appear in multiple domains, so per-domain Term Counts sum to more than the 1239 distinct domain-assigned terms. Total Freq counts all occurrences across the corpus for domain-assigned terms. Bridging Terms indicate terms that co-occur across multiple domain vocabularies. Full per-domain breakdowns are in \texttt{output/data/domain\_statistics.json}.}
\label{tab:terminology_extraction}
\end{table}

Of 11616 total extracted candidate terms, 1239 receive domain
assignments. The global corpus vocabulary possesses a Type-Token Ratio
(TTR) of \textbf{0.0472}, reflecting the dense, highly specialized
nature of the discourse. The absolute highest frequency terms across all
contexts empirically anchor the investigation: \textbf{ant} (20538
occurrences), \textbf{specie} (10623 occurrences), and \textbf{colony}
(7456 occurrences) dominate the conceptual landscape.

Among domains, Behavior \& Identity possesses the highest absolute
occurrence frequency (9239 total occurrences), while Power \& Labor
exhibits the most extensive bridging capacity (139 bridging terms).
Conversely, Economics maintains the most tightly constrained vocabulary
(65 terms) with zero bridging bleed-over (3 bridging terms), reflecting
strict, insular deployment of economic metaphors.

\subsection{Terminology Network
Structure}\label{terminology-network-structure}

Terminology networks were constructed from the term--domain assignments
produced by extraction: two terms are linked when they share at least
one Ento-Linguistic domain, and edge weights are normalized by
domain-set size to emphasize meaningful relationships. Writing \(D(t)\)
for the set of domains term \(t\) is assigned to, an edge is retained
only when its weight exceeds \(0.1\):

\begin{equation}\label{eq:network_edge_weight}
w(u,v) = \frac{|D(u) \cap D(v)|}{\max(|D(u)|, |D(v)|)}
\end{equation}

Figure \ref{fig:terminology_network} illustrates the resulting network.

\begin{figure}[h]
\centering
\includegraphics[width=0.95\textwidth]{/Volumes/external_drive/Git/projects/ongoing/docxology/ento_linguistics/output/figures/terminology_network.png}
\caption{Co-occurrence network of the domain-assigned terminology extracted from the headline abstract layer (7608 open-access PubMed abstracts; 11616 candidate terms, of which 1239 are assigned to at least one Ento-Linguistic domain). Each node is an extracted term: node size is proportional to corpus frequency, and node color encodes the term's primary Ento-Linguistic domain (legend at right). Edge width is proportional to the pipeline relationship weight between terms (shared-domain overlap, Eq.~\ref{eq:network_edge_weight}); isolated terms are omitted, and only the twenty highest-frequency terms are labelled for legibility. Clustered regions are terminology communities dominated by particular domains; terms linking communities are the bridging terms discussed in the text. Network values resolve at build time from \texttt{output/data/domain\_statistics.json} and \texttt{output/data/concept\_map\_summary.json}.}
\label{fig:terminology_network}
\end{figure}

The network exhibits strong modularity: 11622 nodes (11616 extracted
terms plus the 6 conceptual cluster nodes) connected by 1165 edges, with
a clustering coefficient of 0.0413 and average degree of 0.20. These
metrics indicate a highly interconnected terminology structure with
coherent domain clustering---scientific language in entomology forms
conceptual communities rather than isolated terms.

Domain-level network analysis reveals distinct architectures across the
six core themes. As visualized in the aggregate network topology, dense
identity clusters characterize Behavior \& Identity terminology, while
Power \& Labor terminology forms hierarchical, chain-like structures.
Conversely, Sex \& Reproduction terms tend to organize into rigid binary
oppositions, and Economics terms cluster tightly around transactional
frameworks with few bridges to biological mechanism descriptions.

The conceptual bridges between these domains are quantified and
visualized in Figure \ref{fig:domain_overlap}.

\begin{figure}[h]
\centering
\includegraphics[width=0.9\textwidth]{/Volumes/external_drive/Git/projects/ongoing/docxology/ento_linguistics/output/figures/domain_overlap_heatmap.png}
\caption{Cross-domain terminology overlap among the six Ento-Linguistic domains, computed at build time from the term--domain assignments of the headline abstract layer (7608 open-access PubMed abstracts; \texttt{output/data/domain\_statistics.json}). Each cell is the Szymkiewicz--Simpson overlap coefficient (Eq.~\ref{eq:overlap_coefficient}): the number of extracted terms assigned to both domains of the pair, divided by the smaller domain's term count. The heatmap is symmetric, so off-diagonal cells read the same in both directions; darker cells (\texttt{YlOrRd} colormap) indicate higher shared terminology, and the diagonal is 100\% by construction. The strongest observed connectivity joins Power \& Labor with Behavior \& Identity and with Sex \& Reproduction; zero cells (for example Economics with Behavior \& Identity) are observed zeros of the current corpus: no extracted term is assigned to both domains of those pairs.}
\label{fig:domain_overlap}
\end{figure}

Distinctive cross-domain bridges include:

\begin{itemize}
\tightlist
\item
  \textbf{Power \& Labor \(\leftrightarrow\) Behavior \& Identity}:
  Mechanisms of role assignment.
\item
  \textbf{Unit of Individuality \(\leftrightarrow\) Kin \& Relatedness}:
  Foundations of social structure.
\item
  \textbf{Economics \(\leftrightarrow\) Power \& Labor}: Resource
  distribution hierarchies.
\end{itemize}

Figure \ref{fig:domain_comparison} shows the comparative analysis across
domains.

\begin{figure}[h]
\centering
\includegraphics[width=0.9\textwidth]{/Volumes/external_drive/Git/projects/ongoing/docxology/ento_linguistics/output/figures/domain_comparison.png}
\caption{Comparison of terminology characteristics across the six Ento-Linguistic domains, computed at build time from the domain-assigned terminology of the headline abstract layer (7608 open-access PubMed abstracts; \texttt{output/data/domain\_statistics.json}). The 3\,×\,2 grid of bar charts shows (top-left) the number of distinct terms extracted per domain, (top-right) the mean extraction confidence assigned by the term extractor, (center-left) the total corpus frequency of each domain's terms, (center-right) the mean semantic entropy $H(t)$ per domain (Eq.~\ref{eq:semantic_entropy}, averaged over the domain's terms), (bottom-left) the number of bridging terms (terms assigned to more than one domain), and (bottom-right) the mean CACE aggregate score (Clarity, Appropriateness, Consistency, Evolvability; computed per domain on a sample of up to fifty terms). Within each panel, bar height encodes the quantity named on the $y$-axis, with values annotated at the bar tips; bar colors distinguish domains, not magnitudes. Domains with higher mean semantic entropy have usage contexts spanning more sense clusters in this corpus---an observed association, not evidence of causal framing effects.}
\label{fig:domain_comparison}
\end{figure}

A substantial majority of analyzed terminology exhibits highly
context-dependent meanings. Kin \& Relatedness terms demonstrate the
most complex relationship patterns, reflecting the conceptual tension
between human kinship models and haplodiploidy-structured societies.
Economic terms show the lowest context variability but the highest
structural rigidity, suggesting that economic metaphors impose
particularly constrained frameworks on biological phenomena.

\subsection{Framing Analysis}\label{framing-analysis}

Computational identification of framing assumptions reveals systematic
biases embedded within the literature. Anthropomorphic framing
profoundly affects all domains, while hierarchical framing concentrates
heavily within the Power/Labor and Unit of Individuality discourse.

\begin{figure}[h]
\centering
\includegraphics[width=0.9\textwidth]{/Volumes/external_drive/Git/projects/ongoing/docxology/ento_linguistics/output/figures/anthropomorphic_framing.png}
\caption{Inventory of the curated anthropomorphic vocabulary used by the framing analysis over the headline abstract layer (7608 open-access PubMed abstracts). \emph{Left}: number of curated marker terms per category---Hierarchical Terms, Economic Metaphors, Kinship Language, Identity Labels, and Agency Attribution; bar height encodes the count, annotated at the bar tip, with the overall total shown bottom-right. \emph{Right}: up to five example terms per category. These categories define the marker vocabularies that the pipeline's framing analysis matches in term-occurrence contexts; the counts are the sizes of the curated vocabularies, not corpus frequencies. Per-domain anthropomorphic-framing proportions derived from these vocabularies resolve at build time from \texttt{output/data/domain\_statistics.json}.}
\label{fig:anthropomorphic}
\end{figure}

Our ambiguity detection algorithm classifies four distinct ambiguity
types---lexical, contextual, scale-dependent, and temporal---and
confirms that \emph{scale ambiguity} (where meaning shifts across
biological levels of organization) and \emph{context-dependent semantic
drift} are the most prevalent patterns across the corpus (see Section
\ref{sec:supplemental_analysis} for the formal multi-level ambiguity
classification).

\newpage

\section{Results: Domain-Specific Findings}\label{sec:domain_findings}

\subsection{Unit of Individuality}\label{unit-of-individuality}

Frequency and ambiguity analyses show that the highest-frequency terms
(\texttt{colony,\textquotesingle{}\textquotesingle{}}individual'\,')
anchor the domain. Its mean semantic entropy (0.28 bits) sits in the
lower half of the six-domain range (Table
\ref{tab:entropy_distribution}). Figure
\ref{fig:unit_individuality_patterns} details the scale-dependent
terminology patterns within this domain, while per-domain top-term
frequency distributions and word-formation composition breakdowns for
all six domains are visualized in Figure \ref{fig:domain_overview_grid}
and Figure \ref{fig:domain_patterns_grid} respectively.

\begin{figure}[h]
\centering
\includegraphics[width=0.95\textwidth]{/Volumes/external_drive/Git/projects/ongoing/docxology/ento_linguistics/output/figures/domain_overview_grid.png}
\caption{Top terms by corpus frequency for each of the six Ento-Linguistic domains, as a 3\,×\,2 grid of horizontal bar charts computed at build time from the terminology of the headline abstract layer (7608 open-access PubMed abstracts; 1239 domain-assigned terms; \texttt{output/data/domain\_statistics.json}). Within each panel, the ten highest-frequency terms of that domain are ranked by corpus frequency (bar length, annotated at the bar tip); panel titles give each domain's total term count. Bar color encodes per-term semantic entropy $H(t)$ in bits (Eq.~\ref{eq:semantic_entropy}) on a shared \texttt{YlOrRd} scale spanning zero to the corpus maximum, with the shared color bar at right; darker bars indicate terms whose usage contexts span more sense clusters. The grid shows Economics as the smallest domain while its terms carry the highest entropies (Table \ref{tab:entropy_distribution}), with Power \& Labor and Behavior \& Identity also containing high-entropy terms.}
\label{fig:domain_overview_grid}
\end{figure}

\begin{figure}[h]
\centering
\includegraphics[width=0.95\textwidth]{/Volumes/external_drive/Git/projects/ongoing/docxology/ento_linguistics/output/figures/domain_patterns_grid.png}
\caption{Word-formation composition of each Ento-Linguistic domain's vocabulary, shown as donut charts in a 3\,×\,2 grid (one panel per domain), computed at build time from the surface morphology of the domain-assigned terms extracted from the headline abstract layer (7608 open-access PubMed abstracts; \texttt{output/data/domain\_statistics.json}). Each slice is a word-formation class---\emph{hyphenated compound} (e.g. ``queen-worker''), \emph{multiword phrase} (e.g. ``division of labor''), or \emph{single word} (e.g. ``caste'')---with slice angle proportional to that class's share of term counts within the domain. Categories beyond the six largest per domain are grouped as \emph{Other}, and the annotation at each donut's centre gives the domain's term count. Domains whose slices concentrate in hyphenated-compound categories (for example Unit of Individuality) have vocabularies structurally biased toward reified compound labels in this corpus.}
\label{fig:domain_patterns_grid}
\end{figure}

\begin{figure}[h]
\centering
\includegraphics[width=0.9\textwidth]{/Volumes/external_drive/Git/projects/ongoing/docxology/ento_linguistics/output/figures/unit_of_individuality_patterns.png}
\caption{Terminology patterns in the Unit of Individuality domain, computed at build time from the domain's terms in the headline abstract layer (7608 open-access PubMed abstracts; \texttt{output/data/domain\_statistics.json}). \emph{Left}: term-formation patterns (the surface word-formation structure of the domain's terms) as a pie chart, with slice angles proportional to pattern counts and percentages annotated. \emph{Right}: number of domain terms whose names match keyword groups for five biological scales---Colony (Superorganism), Sub-colony (Caste), Individual (Worker), Genomic (Gene-level), and Emergent (Collective)---as a bar chart with counts annotated; counts are floored at one for display, so a near-zero scale is indistinguishable from a single match. Colony-level terms such as ``colony'' and ``superorganism'' populate the colony-scale group; the distribution across scales grounds the scale ambiguities discussed in the text.}
\label{fig:unit_individuality_patterns}
\end{figure}

\subsection{Power \& Labor}\label{power-labor}

The most structurally rigid domain shows clear hierarchical patterns
derived from human social systems
\cite{herbers2007, boomsma2018superorganismality}. Recent molecular
approaches to caste \cite{heinze2017molecular} and epigenetic evidence
that caste determination is a labile developmental process
\cite{warner2024caste} further underscore the need for reform. 75.1\% of
Power \& Labor terms score above baseline on the pipeline's
anthropomorphic-framing proportion for this domain (see
\texttt{domain\_statistics.json}), consistent with pervasive
hierarchical metaphor. Figure \ref{fig:concept_hierarchy} visualizes the
resulting conceptual hierarchy; Figures
\ref{fig:power_labor_frequencies} and \ref{fig:power_labor_ambiguities}
profile the domain's term frequencies and semantic-entropy distribution.
``Caste'' and ``queen'' form central hub terms with the highest
betweenness centrality; ``worker'' and ``slave'' show parasitic
terminology influence \cite{herbers2006}. The chain-like network
structure reflects the linear hierarchies assumed by this vocabulary
rather than the distributed organization documented in behavioral
studies (Figures \ref{fig:concept_hierarchy},
\ref{fig:power_labor_frequencies}, \ref{fig:power_labor_ambiguities}).

\begin{figure}[h]
\centering
\includegraphics[width=0.9\textwidth]{/Volumes/external_drive/Git/projects/ongoing/docxology/ento_linguistics/output/figures/concept_hierarchy.png}
\caption{Concept-centrality structure of the Ento-Linguistic concept map, computed at build time from the concept map built over the headline abstract layer (7608 open-access PubMed abstracts; \texttt{output/data/concept\_map\_summary.json}). Concepts are clusters of domain-assigned terms; a concept's centrality is its number of direct links to other concepts in the map. \emph{Left}: all concepts ranked by centrality (bar length, score annotated at the bar tip), colored green for core concepts (centrality strictly above the map-wide mean) and red for peripheral concepts (at or below the mean). \emph{Right}: centrality against the number of associated terms, with point area proportional to centrality and the ten highest-centrality concepts labelled. The ranking spans all six Ento-Linguistic domains; it is not restricted to Power \& Labor.}
\label{fig:concept_hierarchy}
\end{figure}

\begin{figure}[h]
\centering
\includegraphics[width=0.9\textwidth]{/Volumes/external_drive/Git/projects/ongoing/docxology/ento_linguistics/output/figures/power_and_labor_term_frequencies.png}
\caption{Corpus frequency of the fifteen most frequent Power \& Labor terms, computed at build time from the domain-assigned terminology of the headline abstract layer (7608 open-access PubMed abstracts; \texttt{output/data/domain\_statistics.json}). Bar height encodes total corpus frequency, annotated at the bar tip; bars are ordered by descending frequency, and the \texttt{YlOrRd} shading tracks rank order only, encoding no additional quantity. In the current corpus ``queen,'' ``worker,'' and ``caste'' are the three most frequent terms of this domain, consistent with the entrenched hierarchical framing discussed in the text; the ranking is descriptive, and no significance testing is applied to it.}
\label{fig:power_labor_frequencies}
\end{figure}

\begin{figure}[h]
\centering
\includegraphics[width=0.9\textwidth]{/Volumes/external_drive/Git/projects/ongoing/docxology/ento_linguistics/output/figures/power_and_labor_ambiguities.png}
\caption{Semantic entropy $H(t)$ of Power \& Labor terms (Eq.~\ref{eq:semantic_entropy}), computed at build time from the usage contexts of the domain's terms in the headline abstract layer (7608 open-access PubMed abstracts; TF-IDF vectorization of each term's contexts followed by $k$-means sense clustering; \texttt{output/data/domain\_statistics.json}). The fifteen highest-entropy terms of the domain are shown. \emph{Left}: per-term entropy as bars sorted by descending $H(t)$, each annotated with its entropy in bits and extracted-context count, with a dashed vertical line marking the panel median. \emph{Right}: corpus frequency (horizontal) against $H(t)$ (vertical), with point area proportional to the number of extracted contexts and point color repeating the entropy scale (\texttt{Purples}). ``Caste'' and ``queen'' exhibit elevated entropy in this corpus, consistent with their documented polysemy across hierarchical, reproductive, and behavioral research contexts; entropy values describe context diversity and imply no preferred reform direction.}
\label{fig:power_labor_ambiguities}
\end{figure}

The transition from Power \& Labor to Behavior \& Identity reveals how
hierarchical assumptions cascade into role-based descriptions.

\subsection{Behavior \& Identity}\label{behavior-identity}

Behavioral descriptions create categorical identities that may obscure
the biological fluidity documented in ant task-switching research
\cite{ravary2007, gordon2010}. As \citet{gordon1992wittgenstein}
argues---drawing on Wittgenstein's analysis of category boundaries---the
act of classifying a nestmate as a ``forager'' or a ``nurse'' is not a
neutral observation but an imposition of discrete categories onto
continuous behavioral variation. Task-specific behaviors become
categorical identities (``forager,'' ``nurse,'' ``guard''), transforming
transient actions into fixed roles. Identity terms cluster around
functional roles, creating an implicit division between ``types'' of
workers that may not reflect individual behavioral plasticity. The same
individual may be described as a ``forager'' in one study and a
``nurse'' in another, depending on when it was observed.
\citeauthor{gordon2023ecology}'s \citeyearpar{gordon2023ecology} recent
synthesis demonstrates that task allocation in harvester ant colonies
operates entirely through local interaction networks---brief antennal
contacts modulated by cuticular hydrocarbon profiles---without any
centralized assignment. Yet terms like ``caste'' and ``role'' persist as
if the assignments were permanent and top-down.

Detailed frequency and ambiguity analyses for this domain confirm the
pattern: task-identity terms such as
\texttt{forager\textquotesingle{}\textquotesingle{}\ and}nurse'\,'
exhibit high frequency but moderate-to-high ambiguity (mean semantic
entropy: 0.21 bits), reflecting the gap between categorical labels and
fluid biological reality. Per-domain breakdowns are shown in Figures
\ref{fig:domain_overview_grid} and \ref{fig:domain_patterns_grid}.

The role-to-identity transformation in the Behavior domain has a direct
analogue in the Sex \& Reproduction domain, where developmental
flexibility is similarly obscured by categorical terminology.

\subsection{Sex \& Reproduction}\label{sex-reproduction}

Sex and reproduction terminology shows the lowest overall ambiguity but
reveals a distinctive pattern of \textbf{binary opposition}---the
dominant network structure in this domain (Figure
\ref{fig:terminology_network}). Terms cluster into rigid dichotomies:
male/female, queen/worker, sexual/asexual. These oppositions import
mammalian sex-determination frameworks into a fundamentally different
system: under haplodiploidy, males develop from unfertilized (haploid)
eggs and females from fertilized (diploid) eggs, decoupling sex
determination from the chromosomal mechanisms assumed by standard
terminology \cite{chandra2021epigenetics}. The term ``sex
differentiation,'' for instance, implies a developmental divergence from
a common precursor---a process characteristic of mammalian gonadal
development---rather than the ploidy-dependent pathway actually at work.
Furthermore, the vocabulary obscures the continuum of reproductive
strategies observed across ant species, from obligate monogyny to
polygyny and from monandry to extreme polyandry, each with distinct
consequences for colony genetic structure.

Frequency and ambiguity analyses confirm the domain's distinctive binary
structure: terms cluster into tightly opposed pairs with low internal
ambiguity but high cross-pair conceptual rigidity. Full per-domain
frequency and word-formation patterns are shown in Figures
\ref{fig:domain_overview_grid} and \ref{fig:domain_patterns_grid}.

\subsection{Kin \& Relatedness}\label{kin-relatedness}

Kin and Relatedness terminology exhibits the lowest mean semantic
entropy (0.25 bits) and a web-like network architecture reflecting the
complex, non-intuitive relatedness structures of haplodiploid societies
(Figure \ref{fig:domain_overview_grid}). The central tension is between
human bilateral kinship models---where siblings share \(r = 0.5\)---and
the haplodiploidy-specific asymmetry where full sisters share
\(r = 0.75\) but sisters relate to brothers at only \(r = 0.25\). When
researchers describe colony members as ``sisters,'' the term imports an
assumption of symmetry that masks the very asymmetry on which inclusive
fitness theory depends.

Hub terms such as
\texttt{kin,\textquotesingle{}\textquotesingle{}}relatedness,'\,' and
\texttt{inclusive\ fitness\textquotesingle{}\textquotesingle{}\ bridge\ multiple\ sub-domains.\ Network\ analysis\ reveals\ that\ Hamilton\textquotesingle{}s-rule-adjacent\ vocabulary\ dominates\ the\ discourse,\ often\ at\ the\ expense\ of\ alternative\ frameworks\ such\ as\ multilevel\ selection.\ Analysis\ of\ kinship\ terminology\ shows\ that}kin
selection'\,' co-occurs with
\texttt{altruism\textquotesingle{}\textquotesingle{}\ and}cooperation'\,'
far more frequently than with
\texttt{conflict\textquotesingle{}\textquotesingle{}\ or}policing,'\,'
suggesting a framing bias toward cooperative explanations that may
underrepresent intra-colony conflict dynamics. Per-domain frequency and
pattern breakdowns are provided in Figures
\ref{fig:domain_overview_grid} and \ref{fig:domain_patterns_grid}.

\subsection{Economics}\label{economics}

The Economics domain contains the smallest vocabulary (65 terms) and
zero bridging terms (3)---the most insular domain by a substantial
margin. For comparison, Power \& Labor contributes 139 bridging terms to
adjacent domains, whereas Economics shares vocabulary with none. The
complete term inventory reveals the character of this insularity:
\textbf{allocation} (221 occurrences), \textbf{investment} (277),
\textbf{resources} (712), \textbf{resource} (579), and additional
low-frequency terms including \textbf{trade-off}, \textbf{trade-offs},
\textbf{jack-of-all-trades}, and \textbf{gamma-distribution}. The final
two are anomalies---terms pattern-matched to economics seed vocabulary
that are, in practice, statistical and ecological constructs co-opted by
economic framing, yet their presence reflects how pervasively the
economic paradigm has colonized foraging ecology's conceptual substrate.

The semantically active core terms conflate two fundamentally different
levels of explanation. ``Cost'' may refer to proximate energetic
expenditure (measurable in joules) or to ultimate fitness reduction
(requiring population-level inference); these distinct meanings are
routinely treated as interchangeable. The same proximate--ultimate
conflation operates across ``investment,'' ``resource allocation,'' and
``trade-off.'' The resulting network architecture is self-contained:
transaction-like term pairs (``cost--benefit,''
``allocation--resource'') form tight clusters with 3 bridging edges to
biological-mechanism clusters---indicating that economic terminology
operates as a closed conceptual subsystem that resists integration with
process-level descriptions.

Notably, Economics terms exhibit the highest mean semantic entropy (0.28
bits) of all domains despite zero bridging terms, confirming that
economic metaphors form a self-contained but highly polysemous
subsystem. The average extraction confidence is also the highest,
indicating stable deployment within this insular vocabulary. This
monoculture trades explanatory integration across domains for internal
semantic precision. These patterns are shown across all domains in
Figures \ref{fig:domain_overview_grid} and
\ref{fig:domain_patterns_grid}.

\subsection{Longitudinal Case Studies}\label{longitudinal-case-studies}

To understand how these linguistic paradigms evolve over time, we
conducted longitudinal analysis on two critical terminology clusters:
\emph{Caste} and \emph{Superorganism}.

\subsubsection{Caste Terminology
Evolution}\label{caste-terminology-evolution}

A clear historical trajectory emerges from rigid categories to
plasticity-aware descriptions:

\begin{itemize}
\tightlist
\item
  \textbf{Foundational Era} (pre-1990): Rigid caste categories dominated
  descriptions of task allocation. Terms like ``caste'' and ``subcaste''
  were used as if they denoted fixed, heritable phenotypes---analogous
  to social strata in human societies.
\item
  \textbf{Transitional Era} (1990--2010): Gradual shift toward
  task-based understanding and behavioral ecology.
  \citeauthor{gordon1992wittgenstein}'s
  \citeyearpar{gordon1992wittgenstein} Wittgensteinian critique and
  accumulating behavioral data on task-switching challenged the
  categorical rigidity of caste vocabulary.
\item
  \textbf{Modern Era} (2010--present): Increasing recognition of
  individual variation, behavioral plasticity, and distributed control.
  Epigenetic evidence \cite{chandra2021epigenetics, warner2024caste}
  reveals caste determination as a labile developmental process, further
  undermining the linguistic prior of fixed social categories.
\end{itemize}

\subsubsection{Superorganism Debate: Conceptual
Evolution}\label{superorganism-debate-conceptual-evolution}

The superorganism concept has undergone a parallel structural
transformation. \citeauthor{wheeler1911}'s \citeyearpar{wheeler1911}
early metaphor of the colony-as-organism organized a century of research
while simultaneously constraining how individuality was conceptualized
in social insect biology. More recent work has progressively replaced
this metaphorical convenience with mathematically rigorous, multi-scale
frameworks of biological individuality---particularly the Markov Blanket
formalism \cite{friston2013life} and
\citeauthor{boomsma2018superorganismality}'s
\citeyearpar{boomsma2018superorganismality} analysis of how
``superorganismality'' was lost in translation between evolutionary and
organismic biology. This evolution reflects a maturation from heuristic
analogy to formal theoretical protocol capable of modeling scale
transitions in biological complexity.

\newpage

\section{Discussion}\label{sec:discussion}

\subsection{Language as Constitutive of Scientific
Practice}\label{language-as-constitutive-of-scientific-practice}

Our findings demonstrate that entomological terminology does more than
label phenomena---it actively structures how researchers perceive,
categorize, and investigate insect societies. This result extends the
constructivist tradition in philosophy of science
\cite{latour1987, longino1990} into a domain where the entanglement of
human social concepts with biological description is especially acute.

Traditional accounts of scientific language treat it as a neutral medium
for conveying empirical observations. Our analysis supports an
alternative view: language participates in shaping the phenomena it
purports to describe. When terms such as ``queen'' and ``worker'' are
used to characterize ant colony roles, they import assumptions about
authority, subordination, and fixed identity that may not reflect the
underlying biological organization \cite{herbers2007}. The quantitative
evidence presented in Sections \ref{sec:experimental_results} and
\ref{sec:domain_findings}---particularly the elevated semantic entropy
of Power \& Labor terms and the chain-like network topology that mirrors
human hierarchies rather than stigmergic architectures---provides
corpus-scale empirical support for what previous qualitative critiques
could only assert.

Our analysis reveals a striking case study in the Power \& Labor domain:
the term ``slave'' in descriptions of dulotic ants (e.g.,
\emph{Polyergus} and \emph{Formica sanguinea}). This term, introduced
through early English translations of Pierre Huber's 1810
\emph{Recherches sur les mœurs des fourmis} \cite{huber1810}, carries
deep associations with racialized chattel slavery that reach far beyond
neutral scientific description. More critically, the framing \textbf{may
have discouraged investigation into host resistance} for decades. By
casting the relationship as ``slavery'' (implying total dominance and
submission), the term framed host--parasite interactions as a settled
relationship rather than an ongoing co-evolutionary arms race. Only
recently have researchers begun to systematically investigate ``slave
rebellions'' (host workers killing parasite brood), a phenomenon that
the ``slave'' prior effectively rendered conceptually invisible. Despite
\citeauthor{herbers2006}'s \citeyearpar{herbers2006, herbers2007}
proposed alternatives (``pirate ants'' for the raiders, ``leistic'' for
the behavior), adoption has been slow. Growing professional consensus
within myrmecology acknowledges that reform in entomological terminology
remains overdue, yet institutional inertia and the argument from
literature continuity continue to delay replacement. The Entomological
Society of America's Better Common Names Project
\cite{betternamesproject2024} represents one institutional pathway
forward, but the pace of adoption underscores the depth of
terminological entrenchment analyzed throughout this paper.

This constructive role of language operates at several levels.

At the level of \emph{conceptual framing}, terms carry implicit
theoretical commitments that guide research directions. Our framing
analysis definitively shows that anthropomorphic framing pervades across
all domains, with overt hierarchical framing strongly concentrating in
the Power \& Labor and Unit of Individuality domains. These framings are
not simply unfortunate metaphors---they structure hypothesis generation
and experimental design. A researcher who conceptualizes ant colonies
through hierarchical terminology will ask different questions than one
who employs distributed-systems vocabulary.

At the level of \emph{cross-domain transfer}, terminology borrowed from
human social organization creates systematic biases in how biological
phenomena are interpreted. The chain-like network architecture of Power
\& Labor terminology (Figure \ref{fig:terminology_network}) mirrors the
linear hierarchies of human institutions rather than the distributed,
flexible patterns that behavioral data reveal
\cite{ravary2007, gordon2010}. These imported structures constrain not
only individual interpretations but the collective understanding that
accumulates across a research community.

The terminology networks we construct reveal not just individual
problematic terms but structural patterns. The high clustering
coefficient (0.0413) indicates that terms reinforce each other within
conceptual clusters, creating self-sustaining frameworks that resist
piecemeal reform. This network-level effect connects to
\citeauthor{foucault1972archaeology}'s
\citeyearpar{foucault1972archaeology} analysis of how discursive
formations constrain what can be said and thought within a field, and
extends \citeauthor{lakoff1980metaphors}'s
\citeyearpar{lakoff1980metaphors} demonstration of pervasive
metaphorical reasoning into formal scientific discourse. Moreover, as
recent accounts of collective behavior
\cite{gordon2019ecology, gordon2023ecology} gain traction, the need for
precise language to distinguish between metaphorical mapping and
functional identity becomes even more critical.

\subsection{From Metaphor to Mechanism: An Active Inference
Perspective}\label{from-metaphor-to-mechanism-an-active-inference-perspective}

Viewing ant colonies through an Active Inference lens
\cite{friston2010free, clark2013whatever} fundamentally reframes the
relationship between language and scientific understanding. Under this
framework, terminology constitutes the \textbf{prior beliefs} of a
generative model. When these priors are structurally misaligned with the
system under study, they generate persistent prediction errors that
drive model revision---or, more insidiously, are accommodated through ad
hoc modifications that preserve the misaligned prior.

The Active Inferants framework \citep{friedman2021active} makes this
tension especially vivid. \citet{friedman2021active} demonstrate that
ant colonies can be modeled as ensembles of active inference
agents---each individual performing approximate Bayesian inference over
local pheromone gradients---whose collective behavior emerges from
stigmergic coupling without any centralized controller. This model
succeeds precisely \emph{because} it abandons the monarch-and-subject
vocabulary embedded in traditional terminology. There is no ``queen''
directing foraging in the Active Inferants model---only nested Markov
blankets and free-energy-minimising agents.

This perspective aligns with what we term \textbf{Environment-Centric
Active Inference (EC-AIF)}---a synthesis drawing on niche-construction
and active inference principles---which defines an ``individual'' not by
its skin but by its \emph{niche}---the set of states it can
statistically regulate. In EC-AIF, the ``individual'' ant and the
``colony'' superorganism are simply two different scales of niche
construction \cite{deacon2011incomplete}. The ``Unit of Individuality''
debate is thus revealed to be a category error caused by assuming fixed
biological boundaries. Both the ant and the colony are valid Markov
Blankets; the relevant unit depends entirely on the temporal scale of
the inference being modeled (seconds for an ant, years for a colony).

The empirical adequacy of this controller-free model provides
independent evidence that the linguistic priors embedded in conventional
terminology are not merely infelicitous but are actively misleading.

In the \textbf{Free Energy Principle} framework, biological systems
maintain their integrity by minimizing variational free
energy---essentially, by acting to fulfill the predictions of their
generative models \cite{friston2013life}.

When researchers model these systems using hierarchical language
(``queen control''), they impose a scientific generative model that
assumes \textbf{centralized prediction-error minimization}. However, ant
colonies exist through \textbf{distributed active inference}: each
individual acts on local Markovian states (pheromones, tactile cues)
without a global representation of the colony state.

By misidentifying the \textbf{locus of agency}---attributing it to a
``queen'' rather than the collective manifold---scientific terminology
introduces a formal \textbf{modeling error}. This error forces
researchers to postulate ad hoc mechanisms (such as ``police'' workers
or ``royal decrees'') to explain deviations from the hierarchical prior.
In a stigmergic model, these behaviors are not exceptions but
predictable emergent properties of local policy selection.

A concrete example clarifies the stakes. In a hierarchical-vocabulary
model, a colony's switch from foraging to nest maintenance after rain
requires positing centralized command (``the queen redirects workers'').
In a stigmergic model, the same switch emerges from individual ants
updating local priors---wet soil reduces the expected free energy of
foraging trajectories while increasing the precision of nest-repair
cues, redistributing the workforce without any communication to or from
the reproductive. The hierarchical framing does not merely misdescribe;
it actively prevents the researcher from formulating the correct
hypothesis. Terminology reform is therefore a process of \textbf{model
selection}: replacing high-entropy priors with lower-entropy,
mechanistically accurate descriptors.

\subsection{Comparison with Existing
Approaches}\label{comparison-with-existing-approaches}

Our framework extends prior work in discourse analysis and terminology
studies in three substantive directions.

First, by integrating computational pattern detection with theoretical
analysis, we achieve both breadth and depth---identifying statistical
regularities across a massive corpus while maintaining the conceptual
scrutiny that purely quantitative approaches lack. Existing
computational approaches to scientific discourse
\cite{chen2006citespace} primarily model citation networks rather than
the semantic content of terminological usage. Qualitative critiques of
loaded scientific language \cite{herbers2007, keller1991language} offer
incisive analysis of individual terms but cannot capture systemic
patterns. Our framework bridges this gap, supporting both SSK arguments
about social construction of scientific facts \cite{latour1987} and
feminist epistemological critiques of androcentric category projection
\cite{haraway1991}.

Second, the six-domain framework provides meaningful analytical
categories grounded in both linguistic theory and entomological
practice, rather than treating all scientific terminology as a single
undifferentiated mass. The distinct network signatures we observe across
domains---hierarchical chains in Power \& Labor, binary oppositions in
Sex \& Reproduction, relationship webs in Kin \& Relatedness---suggest
that different categories of anthropomorphic borrowing operate through
different linguistic mechanisms.

Third, the CACE meta-standards (Section \ref{sec:methodology}) offer a
concrete evaluation framework that moves beyond critique toward
constructive reform. Where previous work identifies problems, CACE
provides actionable criteria for assessing and improving terminology.

\subsection{Practical Implications for Scientific
Communication}\label{practical-implications-for-scientific-communication}

\subsubsection{Terminology Awareness and
Reform}\label{terminology-awareness-and-reform}

Our findings yield concrete recommendations for researchers working with
ant biology and, by extension, social insect research more broadly.

Researchers should become intensely aware of how their terminological
choices import assumptions. The domain contrasts that the statistical
battery supports---the elevated semantic entropy of Economics
terminology and the dense bridging vocabulary of Power \& Labor (139
bridging terms)---trace exactly the contours where linguistic precision
would most improve scientific communication. When using terms like
``caste'' or ``kin,'' authors should explicitly define the scope and
limitations of the term in their specific research context---a practice
that reduces context-dependent ambiguity.

Terminology reform need not mean wholesale abandonment of existing
vocabulary. Instead, we advocate for \emph{qualified usage}: retaining
familiar terms where they are genuinely informative while flagging their
metaphorical status and providing operational definitions. ``Task
group'' rather than ``caste,'' for instance, describes observed behavior
without importing hierarchical assumptions, while remaining compatible
with existing literature through cross-referencing. Recent community
efforts such as the ESA Better Common Names Project
\cite{betternamesproject2024} and \citeauthor{herbers2007}'s
\citeyearpar{herbers2007} call for language reform provide models for
systematic terminology revision.

\subsubsection{Cross-Domain
Communication}\label{cross-domain-communication}

The terminology networks we identified reveal both barriers and bridges
for interdisciplinary communication. Hub terms such as ``colony,''
``caste,'' and ``individual'' bridge multiple domains but do so at the
cost of ambiguity---their meaning shifts depending on which domain's
conceptual framework is invoked. Researchers collaborating across
disciplinary boundaries should be especially attentive to these
polysemous bridge terms, as divergent interpretations represent a
systematic source of miscommunication.

Conversely, the strong domain clustering (clustering coefficient 0.0413)
indicates that within-domain communication is relatively coherent. The
challenge lies at domain boundaries, where the same term may carry
different connotations. Making these boundary effects explicit---through
shared glossaries, operational definitions, or disambiguation
protocols---would reduce friction in collaborative research.

\subsection{The ``Slave'' Terminology Debate: A Case Study in
Reform}\label{the-slave-terminology-debate-a-case-study-in-reform}

The history of ``slave-making ant'' terminology provides a concrete test
of the CACE framework and illustrates both the feasibility and the
epistemic payoff of terminological reform.

For over a century, species such as \emph{Polyergus} and \emph{Formica
sanguinea} were described through a master--slave metaphor: raided brood
were ``slaves,'' raiding species were ``slave-makers,'' and the behavior
itself was ``slave-making'' \cite{hollldobler1990}.
\citet{herbers2006, herbers2007} catalysed reform by demonstrating that
the terminology naturalized a human institution of extreme moral weight
while simultaneously obscuring the biology. Evaluating ``slave'' through
CACE makes the case transparent:

\begin{itemize}
\tightlist
\item
  \textbf{Clarity}: ``Slave'' conflates the social relationship
  (exploited labor under coercion) with the biological mechanism (brood
  parasitism and chemical manipulation of host behavior). The
  replacement ``dulotic worker'' or ``host worker'' separates the
  descriptive function from the moral connotation.
\item
  \textbf{Appropriateness}: Enslaved humans exercise agency, resistance,
  and cultural production; parasitized ant brood do not. The metaphor
  projects attributes absent from the target phenomenon.
\item
  \textbf{Consistency}: ``Slave'' was applied inconsistently---sometimes
  to the individual host worker, sometimes to the entire host colony,
  and occasionally to unrelated phenomena such as facultative social
  parasitism.
\item
  \textbf{Evolvability}: Modern understanding of superorganism-level
  immune responses and chemical mimicry \cite{wilson2008superorganism}
  renders the ``slave'' metaphor actively misleading, since the host
  workers' behavior results from chemical deception rather than
  submission.
\end{itemize}

The shift to ``social parasitism,'' ``dulosis,'' and ``host worker'' in
journals including \emph{Insectes Sociaux} and \emph{Behavioral Ecology}
demonstrates that terminological reform need not sever continuity with
the literature: systematic cross-referencing and the indexing capacity
of modern databases ensure discoverability. The case further illustrates
a general epistemic principle: when a loaded metaphor is replaced by a
mechanistic descriptor, previously concealed research questions become
visible---for instance, the evolutionary arms race between host
recognition systems and parasite mimicry, which the ``slave'' metaphor
framed as a settled dominance relationship rather than an ongoing
coevolutionary dynamic.

Quantitative CACE scoring confirms this qualitative assessment.
Aggregate scores rise from 0.50 (``slave'') to 0.47 (``host worker''),
with Appropriateness increasing from 0.50 to 0.40 (severing the
anthropomorphic linkage eliminates the penalty entirely) and Clarity
from 1.00 to 1.00 (reduced semantic entropy reflecting a mechanistically
specific dulosis descriptor). All values are rendered from the frozen
statistics artifact (\texttt{output/data/statistical\_analysis.json}) at
build time. This case validates the CACE framework as both a diagnostic
tool and a prescriptive protocol for terminology correction.

\subsection{Limitations}\label{limitations}

Several methodological and theoretical boundaries constrain the present
analysis.

\begin{enumerate}
\def\labelenumi{\arabic{enumi}.}
\tightlist
\item
  \textbf{Corpus scope}: Analysis is limited to English-language
  publications; multilingual patterns remain unexplored. Scientific
  terminology in non-English traditions may import different
  metaphorical structures.
\item
  \textbf{Text accessibility}: Full-text availability varies by
  publication date and venue, introducing potential sampling bias toward
  more recent and open-access literature.
\item
  \textbf{Context window size}: Co-occurrence analysis uses configurable
  sliding windows (10-word default for term-level, 50-word for
  domain-level); longer-range conceptual relationships may be missed.
\item
  \textbf{Domain boundaries}: The six Ento-Linguistic domains were
  defined \emph{a priori} from seed lists; some terms (e.g., ``colony'')
  span multiple domains, creating classification challenges. Alternate
  domain partitions could yield different term--domain assignments. Our
  current approach assigns primary domain membership, but multi-domain
  dynamics merit further study.
\item
  \textbf{Historical depth}: Cross-sectional analysis does not fully
  capture the temporal evolution of terminological usage, though our
  case studies (Section \ref{sec:supplemental_analysis}) offer
  preliminary longitudinal evidence.
\item
  \textbf{Interdisciplinary borrowing}: The extent to which
  entomological terminology is shaped by borrowing from economics,
  sociology, and political science is not yet quantified systematically.
\item
  \textbf{Functional heterogeneity}: Some terminology may function
  differently across phases of inquiry---metaphorical during hypothesis
  generation but operationally precise during data collection---a
  dynamic our static analysis cannot fully capture.
\end{enumerate}

Future research directions---including multilingual comparative
analysis, longitudinal corpus studies, and educational applications of
the CACE meta-standards---are developed in Section \ref{sec:conclusion}.

\newpage

\section{Conclusion}\label{sec:conclusion}

This work establishes Ento-Linguistic analysis as a methodology for
examining how scientific language constitutes---rather than merely
represents---knowledge about insect biology. Through computational
analysis of terminology networks across \textbf{7608 entomological
publications} (999905 tokens; 11616 extracted candidate terms, 1239
domain-assigned) and six analytically distinct domains, we demonstrate
that entomological terminology carries systematic patterns of ambiguity,
anthropomorphic framing, and conceptual structure that actively shape
research practice. The accompanying open-source computational pipeline
provides a reproducible toolkit for extending this analysis to new
corpora and domains.

\subsection{Core Contributions}\label{core-contributions}

The work makes three primary contributions. First, the six-domain
analytical framework provides a comprehensive, reproducible architecture
for examining how language shapes scientific understanding in entomology
and, by extension, in other fields where human social concepts are
projected onto non-human systems. Second, the computational pipeline
demonstrates that large-scale, quantitative analysis of scientific
discourse is both feasible and revealing---exposing structural patterns
that qualitative analysis alone cannot detect. Third, the CACE
meta-standards, defined in Section \ref{sec:methodology}, offer a
practical evaluation framework:

\begin{itemize}
\tightlist
\item
  \textbf{Clarity}: stable, non-ambiguous definitions across scales
\item
  \textbf{Appropriateness}: metaphors apt for the biological phenomenon
\item
  \textbf{Consistency}: uniform usage within and across the field
\item
  \textbf{Evolvability}: robustness to new empirical discoveries
\end{itemize}

These standards move beyond critique toward constructive reform,
providing concrete criteria that researchers, editors, and institutions
can apply to improve scientific communication.

The quantitative reach of these findings underscores their significance.
Across the 1239 domain-assigned terms extracted from 7608 publications,
12.8\% exhibit highly context-dependent meanings. The 6 conceptual
clusters identified in the concept map (linked by 15 weighted
relationships) confirm that the terminological landscape is both deeply
interconnected and systematically biased. The Power \& Labor
domain---containing the most entrenched anthropomorphic
vocabulary---generates the strongest cross-domain interference, with 139
bridging terms propagating hierarchical framing into adjacent domains.
The Economics domain, despite its tightly constrained 65-term vocabulary
with 3 bridging terms, exhibits both the highest mean semantic entropy
and the greatest proportion of high-entropy terms, indicating that
economic metaphors form a self-contained but intensely polysemous
subsystem. Crucially, CACE validation on the ``slave'' \(\rightarrow\)
``host worker'' terminological reform demonstrates significant
overarching score improvement, confirming that the framework functions
as both an analytical diagnostic and a prescriptive template for
actionable reform.

\subsection{Future Directions}\label{future-directions}

Several avenues emerge for extending this work.

\textbf{Multilingual and Cross-Cultural Analysis.} Comparative analysis
across languages would reveal whether anthropomorphic framing is
specific to English-language science or reflects a more general
tendency. Preliminary evidence from German (\emph{Königin},
\emph{Arbeiterin}) and Japanese entomological traditions suggests both
convergence and divergence in metaphorical borrowing, warranting
systematic investigation.

\textbf{Longitudinal Terminology Tracking.} Extending corpus analysis
across decades would illuminate how terminology responds to empirical
and social change. Do genomic discoveries erode the dominance of
``caste'' vocabulary? Does institutional reform (e.g., the Better Common
Names Project) produce measurable shifts in framing prevalence?
Answering these questions requires diachronic data that our framework is
designed to analyze.

\textbf{Educational and Editorial Tools.} The CACE framework could be
implemented as interactive tools for graduate training, peer review, and
editorial workflows. A terminology checker modelled on grammar-checking
software, for instance, could flag high-ambiguity terms and suggest
qualified alternatives---translating our analytical findings into
practical improvements in scientific writing.

\textbf{Cross-Disciplinary Extension.} The Ento-Linguistic framework is
not specific to entomology. Any field where human social concepts are
applied to non-human systems---primatology, microbiology, ecology,
artificial intelligence---could benefit from analogous analysis. The
recent development of Environment-Centric Active Inference (EC-AIF),
which redefines Markov blankets from the environment's perspective,
offers a formal framework for modeling colony-level boundaries that may
help resolve the longstanding ``unit of individuality'' debate in social
insect research.

\textbf{Cross-Era Semantic Meta-Analysis.} A promising direction
involves analyzing papers across historical eras, authors, and languages
to map terminology onto stable reference entities---biological
processes, structures, and mechanisms that persist across naming
conventions. By grounding each terminological variant (e.g., ``queen,''
``gyne,'' ``primary reproductive,'' \emph{Königin}) to a shared
ontological referent, comprehensive meta-analysis of scientific
\emph{semantics} becomes possible, not merely syntax.

Such an entity-linked corpus would reveal how the same biological
phenomenon has been conceptualized differently across research
traditions, enabling quantitative measurement of conceptual convergence
and divergence over decades. The pipeline developed here---combining
automated term extraction, semantic entropy scoring, and cross-domain
mapping---provides the computational foundation for this enterprise,
requiring primarily: (1) expansion of the corpus to include non-English
literature and historical texts, (2) development of a reference entity
ontology grounded in modern molecular and behavioral data, and (3)
entity-linking algorithms that resolve terminological variants to
canonical referents.

\subsection{Closing Remarks}\label{closing-remarks}

The entanglement of speech and thought in scientific practice is neither
accidental nor inconsequential. When a researcher describes
\emph{Diacamma} nestmates as ``queens'' and ``workers,'' these terms
carry an entire social ontology that may obscure the fluid,
experience-dependent task performance documented by \citet{ravary2007}.
Replacing ``queen'' with ``primary reproductive'' is not cosmetic---it
is an act of \textbf{model repair}, aligning our linguistic priors with
the physics of distributed systems and reducing the \textbf{variational
free energy} of our scientific explanations.

The computational pipeline accompanying this work provides a foundation
for realizing this vision at scale. Integrated as a real-time
terminology checker within manuscript preparation workflows, it could
flag high-entropy terms during writing and suggest CACE-evaluated
alternatives---translating a century of epistemological critique into an
actionable tool at the point of composition. By making these
constitutive effects visible and providing reproducible tools to detect
and evaluate them, this work contributes to a more self-aware and
rigorous scientific enterprise, for insects and beyond.

\newpage

\section{Related Work}\label{sec:related_work}

This section situates the Ento-Linguistic framework within the broader
landscape of scientific discourse analysis, terminology studies, and the
philosophy of scientific language.

\subsection{Critical Discourse Analysis and Science
Studies}\label{critical-discourse-analysis-and-science-studies}

The tradition of critical discourse analysis (CDA), as formalized by
\citet{fairclough1992} and extended by \citet{wodak2009methods},
provides the methodological foundation for examining how language
structures power relations and institutional knowledge. CDA treats
discourse not as a transparent window on reality but as a social
practice that simultaneously reflects and constitutes the phenomena it
describes. Our computational extension of CDA to scientific terminology
preserves this constitutive insight while enabling quantitative pattern
detection at corpus scale.

Within the sociology of scientific knowledge (SSK), \citet{latour1987}
demonstrated how scientific facts are constructed through networks of
human actors, instruments, and inscriptions---of which terminology is a
central component. \citet{hacking1999social} refined the constructionist
position by distinguishing between the social construction of
\emph{ideas} about natural kinds and the construction of the kinds
themselves, a distinction directly relevant to entomological
terminology: the term ``caste'' constructs a framework for understanding
ant social organization, but the behavioral phenotypes it labels are
empirically real. Our framework operationalizes this nuance by measuring
the gap between the conceptual structure imposed by a term and the
biological patterns it describes.

\citeauthor{kuhn1996}'s \citeyearpar{kuhn1996} analysis of paradigm
shifts highlighted how shared vocabulary both enables and constrains
scientific communities. The terminology networks we construct (Section
\ref{sec:experimental_results}) provide empirical evidence for Kuhnian
incommensurability at the linguistic level: domain-specific vocabulary
clusters resist integration, and paradigm-bridging terms carry high
ambiguity precisely because they must reconcile incompatible conceptual
frameworks. \citeauthor{wheeler1911}'s \citeyearpar{wheeler1911} early
framing of the ant colony as an ``organism'' exemplifies this
process---a metaphor that organized a century of research while
simultaneously constraining how individuality was conceptualized in
social insect biology.

\subsection{Feminist and Postcolonial
Epistemology}\label{feminist-and-postcolonial-epistemology}

Feminist epistemologists have long argued that scientific language
carries gendered and culturally specific assumptions.
\citet{keller1991language} demonstrated how metaphors of mastery and
control pervade biological explanation, and \citet{haraway1991} showed
how primatology's anthropomorphic vocabulary reflects Western gender
norms projected onto non-human societies. \citet{longino1990} argued
that the objectivity of science depends on critical community scrutiny
of precisely the kind of background assumptions that terminology
encodes.

Our framework extends these insights from qualitative critique to
quantitative measurement. The framing prevalence analysis presented in
Section \ref{sec:supplemental_analysis} provides empirical evidence for
the anthropomorphic and hierarchical framings that critics have
identified qualitatively. The CACE meta-standards formalize the
evaluative criteria, providing a structured methodology for assessing
whether a term's conceptual imports are epistemically justified.

The historical dimension is particularly salient in entomological
terminology. Terms like ``slave'' and ``caste'' import specific
historical assumptions about social organization that do not align with
modern biological understanding \cite{herbers2006, herbers2007}.
Historical analysis reveals that early entomology often employed
metaphors of hierarchy and control to describe insect behavior,
influenced by the social contexts of the time
\cite{mavhunga2018transient, sleigh2007ants}. The persistence of these
historical artifacts in modern scientific naming continues to obscure
biological reality, as colonial-era epistemological frameworks remain
embedded in entomological vocabulary \cite{vis2026colony}.
\citeauthor{berlin1992}'s \citeyearpar{berlin1992} cross-cultural
studies of biological classification demonstrate that alternative
taxonomic systems---grounded in different cultural assumptions---are
equally effective for organizing biological knowledge. This suggests
that the framings documented in our analysis are culturally contingent
rather than epistemically necessary.

\subsection{Computational Approaches to Scientific
Discourse}\label{computational-approaches-to-scientific-discourse}

Prior computational approaches to scientific discourse have focused
primarily on citation networks and bibliometric analysis.
\citeauthor{chen2006citespace}'s CiteSpace framework
\citeyearpar{chen2006citespace} maps the intellectual structure of
research fields through co-citation patterns, but does not analyze the
semantic content of terminology. Natural language processing
applications in biomedicine---including biomedical named entity
recognition and terminological relation extraction---optimize for
information extraction rather than conceptual critique.

Our framework occupies a distinct position: it combines the analytical
depth of CDA with the scalability of computational text processing,
targeting the \emph{conceptual implications} of terminology rather than
merely identifying or extracting terms. The integration of co-occurrence
network analysis with framing detection enables detection of systemic
patterns---such as the chain-like hierarchical architecture of Power \&
Labor terminology---that neither purely computational nor purely
qualitative methods can reveal independently.

\subsection{Terminology Studies in
Entomology}\label{terminology-studies-in-entomology}

Within entomology specifically, debates over the adequacy of inherited
terminology have a long genealogy. \citet{wheeler1911} systematized the
organismic metaphor of the colony-as-superorganism, establishing a
vocabulary whose hierarchical assumptions still structure modern
discourse. The philosophical stakes of this vocabulary were first
articulated by \citet{gordon1992wittgenstein}, who applied
Wittgenstein's analysis of category boundaries to argue that the act of
classifying individual ants as ``foragers,'' ``nurses,'' or ``soldiers''
is not a neutral empirical operation but a theory-laden imposition of
discrete categories onto continuous behavioral variation. Gordon's
insight---that the observer's categorical vocabulary determines what
counts as an observation---directly anticipates our computational
measurement of the same phenomenon through semantic entropy: terms whose
usage contexts span many distinct senses (high \(H(t)\)) are precisely
those whose categorical boundaries are most contested.

\citet{herbers2006, herbers2007} initiated the modern institutional
debate over loaded language in social insect research, focusing on
racially charged metaphors such as ``slave raid'' and ``slave-making
ant.'' This critique catalyzed the Entomological Society of America's
Better Common Names Project \cite{betternamesproject2024}, the most
systematic institutional effort at terminological reform to date, which
established formal guidelines for replacing culturally loaded common
names with descriptively accurate alternatives.
\citet{boomsma2018superorganismality} traced how the superorganism
concept was ``lost in translation'' between different theoretical
frameworks---a case study in the terminological dynamics our framework
is designed to detect. \citet{sleigh2007ants} provided a cultural
history of myrmecology that documents how broader social and cultural
currents have shaped the language of ant research across centuries.
Recent epigenetic research further undercuts the biological
justification for rigid ``caste'' terminology: \citet{warner2024caste}
show that caste differentiation in ants becomes increasingly
\emph{canalized} from early development through cascading
gene-expression changes modulated by juvenile hormone signaling---a
fundamentally labile process that the term ``caste'' misleadingly
implies is fixed.

More broadly, the need to broaden conceptions of social insects beyond
the traditional eusociality framework
\cite{boomsma2018superorganismality} implicitly challenges the
terminology built around that framework---particularly ``caste,''
``queen,'' and ``worker'' as universalized descriptors of insect social
organization. Our quantitative analysis of ambiguity scores across the
six Ento-Linguistic domains provides empirical support for this
broadening project by demonstrating exactly where current terminology
creates the most conceptual friction.

\subsection{Active Inference and Colony
Modeling}\label{active-inference-and-colony-modeling}

The Free Energy Principle and Active Inference
\cite{friston2010free, friston2013life} provide the theoretical backbone
for our analysis. \citeauthor{clark2013whatever}'s
\citeyearpar{clark2013whatever} predictive processing framework
establishes the cognitive context in which language acts as a
hyper-prior, and \citeauthor{kirchhoff2018markov}'s
\citeyearpar{kirchhoff2018markov} application of Markov blankets to
biological systems supports our analysis of how terminology
mis-specifies system boundaries.

Most directly relevant is the \emph{Active Inferants} framework of
\citet{friedman2021active}, who model ant colony foraging as a
multiscale ensemble of active inference agents. Each ant performs
approximate Bayesian inference over local pheromone gradients, and
collective behavior emerges through stigmergic coupling---a mechanism
first formalized by \citet{grasse1959reconstruction} as indirect
coordination through environmentally mediated traces---without
centralized control. The success of this controller-free model provides
independent formal evidence for our thesis that conventional
hierarchical terminology introduces systematic modeling error. This
perspective also intersects with the eusociality debate catalyzed by
\citet{nowak2010evolution}, who challenged kin-selection explanations of
eusociality by demonstrating that standard natural-selection models with
population structure suffice---an argument that, regardless of its
contested status, underscores how terminological commitments (e.g.,
\texttt{kin\ selection\textquotesingle{}\textquotesingle{}\ vs.\textbackslash{}}multilevel
selection'\,') frame theoretical controversies. Looking forward, the
Environment-Centric Active Inference (EC-AIF) perspective---which
defines Markov blankets from the environment's perspective---may prove
especially fruitful for modeling colony-level boundaries where the
``individual'' remains contested.

\subsection{Positioning This Work}\label{positioning-this-work}

Our contribution is distinguished from prior work along three axes.
\emph{Methodologically}, we integrate computational and theoretical
approaches in a bidirectional iterative process rather than treating
them as independent tracks. \emph{Analytically}, the six-domain
framework provides a comprehensive yet tractable decomposition of the
problem space, grounded in both linguistic theory and entomological
practice. \emph{Pragmatically}, the CACE meta-standards offer a
constructive evaluation framework that moves beyond critique to provide
actionable criteria for terminological improvement---criteria validated
by the historical case of ``slave'' terminology reform (Section
\ref{sec:discussion}).

\newpage

\section{Acknowledgments}\label{sec:acknowledgments}

We gratefully acknowledge the contributions of individuals and
institutions that made this research possible.

\subsection{Institutional Support}\label{institutional-support}

This work was conducted at the Active Inference Institute. We thank the
Institute for providing the research environment and collaborative
infrastructure that supported the development of the Ento-Linguistic
framework.

\subsection{Collaborations}\label{collaborations}

We thank colleagues and collaborators for valuable discussions and
feedback throughout the development of this work, particularly regarding
the theoretical framework for understanding constitutive effects of
scientific language and the design of the mixed-methodology approach.

\subsection{Data and Software}\label{data-and-software}

This research builds upon open-source software tools and publicly
available datasets. We acknowledge:

\begin{itemize}
\tightlist
\item
  Python scientific computing stack (NumPy, SciPy, Matplotlib, NetworkX)
\item
  Natural Language Toolkit (NLTK) for text processing and scikit-learn
  for validation
\item
  LaTeX and Pandoc for document preparation
\item
  Published entomological literature informing the domain terminology
  seeds
\end{itemize}

\begin{center}\rule{0.5\linewidth}{0.5pt}\end{center}

\emph{All errors and omissions remain the sole responsibility of the
authors.}

\newpage

\section{Symbols and Notation Glossary}\label{sec:glossary}

This glossary defines the mathematical notation and domain-specific
terminology used throughout the manuscript.

\subsection{Mathematical Notation}\label{mathematical-notation}

{\def\LTcaptype{none} % do not increment counter
\begin{longtable}[]{@{}
  >{\raggedright\arraybackslash}p{(\linewidth - 4\tabcolsep) * \real{0.2500}}
  >{\raggedright\arraybackslash}p{(\linewidth - 4\tabcolsep) * \real{0.4062}}
  >{\raggedright\arraybackslash}p{(\linewidth - 4\tabcolsep) * \real{0.3438}}@{}}
\toprule\noalign{}
\begin{minipage}[b]{\linewidth}\raggedright
Symbol
\end{minipage} & \begin{minipage}[b]{\linewidth}\raggedright
Description
\end{minipage} & \begin{minipage}[b]{\linewidth}\raggedright
First Use
\end{minipage} \\
\midrule\noalign{}
\endhead
\bottomrule\noalign{}
\endlastfoot
\(G = (V, E)\) & Terminology network (graph with vertices and edges) &
Eq. \ref{eq:network_edge_weight} \\
\(D(t)\) & Set of Ento-Linguistic domains term \(t\) is assigned to &
Eq. \ref{eq:network_edge_weight} \\
\(w(u,v)\) & Edge weight between terms \(u\) and \(v\): shared-domain
overlap \(|D(u) \cap D(v)| / \max(|D(u)|, |D(v)|)\) & Eq.
\ref{eq:network_edge_weight} \\
\(H(t)\) & Semantic entropy of term \(t\) in bits (Shannon entropy over
usage-context clusters) & Eq. \ref{eq:semantic_entropy} \\
\(H^*\) & High-entropy threshold (2.0 bits, \(\geq 4\) equiprobable
senses) & Eq. \ref{eq:semantic_entropy} \\
\(H_{\max}\) & Maximum attainable entropy for \(k\) clusters,
\(\log_2 k\) & Eq. \ref{eq:semantic_entropy} \\
\(\hat{H}(t)\) & Normalized semantic entropy
\(H(t) / H_{\max} \in [0,1]\) & Eq. \ref{eq:semantic_entropy} \\
\(p_i\) & Empirical proportion of contexts assigned to semantic cluster
\(i\) & Eq. \ref{eq:semantic_entropy} \\
\(k\) & Number of semantic sense clusters (\(k\)-means,
\(k = \max(2,\ \min(k_{\max},\ n{-}1,\ \max(2, \lfloor\!\sqrt{n}\rfloor)))\);
\(k_{\max}=5\), \(n=|C_t|\); \(k < n\)) & Eq.
\ref{eq:semantic_entropy} \\
\(C_t\) & Set of valid usage contexts of term \(t\) (sentences with
\(\geq 3\) words) & Eq. \ref{eq:semantic_entropy} \\
\(S_t\) & Set of biological scale levels expressed in term \(t\)'s
contexts & Eq. \ref{eq:cace_evolvability} \\
\(w_{AB}\) & Overlap coefficient (Szymkiewicz--Simpson) between concept
sets \(A\) and \(B\) & Eq. \ref{eq:overlap_coefficient} \\
\(w_\text{base}\) & Base overlap-coefficient weight in composite
relationship strength & Sec. \ref{sec:methodology} \\
\(r_\text{term}\) & Term-overlap ratio component of composite
relationship strength & Sec. \ref{sec:methodology} \\
\(r_\text{domain}\) & Domain-overlap ratio component of composite
relationship strength & Sec. \ref{sec:methodology} \\
\(\text{Clarity}(t)\) & CACE Clarity score:
\(\max(0, 1 - H(t)/\log_2 10)\) & Eq. \ref{eq:cace_clarity} \\
\(\text{Appropriateness}(t)\) & CACE Appropriateness score (penalizes
anthropomorphic terms) & Eq. \ref{eq:cace_appropriateness} \\
\(\text{Consistency}(t)\) & CACE Consistency score: mean pairwise cosine
similarity of context vectors & Eq. \ref{eq:cace_consistency} \\
\(\text{Evolvability}(t)\) & CACE Evolvability score: proportion of
biological scale levels in contexts & Eq. \ref{eq:cace_evolvability} \\
\(\mathcal{A}\) & Set of anthropomorphic terms (queen, king, slave,
worker, soldier, nurse, \ldots) & Eq. \ref{eq:cace_appropriateness} \\
\(F(D, T)\) & Discursive framing network function for domain \(D\) and
term set \(T\) & Supplemental Eq. \ref{eq:discursive_framing} \\
\(M_{ij}\) & Cross-domain mapping strength between domains \(D_i\) and
\(D_j\) & Supplemental Eq. \ref{eq:cross_domain_mapping} \\
\(\Delta G(t)\) & Temporal network evolution (graph change over time) &
Supplemental Eq. \ref{eq:temporal_network_evolution} \\
\(B\) & Markov Blanket boundary of a system & Supplemental Eq.
\ref{eq:markov_blanket} \\
\(\mu\) & Internal states (conditionally independent of external given
blanket) & Supplemental Eq. \ref{eq:markov_blanket} \\
\(\eta\) & External states & Supplemental Eq. \ref{eq:markov_blanket} \\
\end{longtable}
}

\subsection{Theoretical Terms}\label{theoretical-terms}

{\def\LTcaptype{none} % do not increment counter
\begin{longtable}[]{@{}
  >{\raggedright\arraybackslash}p{(\linewidth - 4\tabcolsep) * \real{0.2222}}
  >{\raggedright\arraybackslash}p{(\linewidth - 4\tabcolsep) * \real{0.4444}}
  >{\raggedright\arraybackslash}p{(\linewidth - 4\tabcolsep) * \real{0.3333}}@{}}
\toprule\noalign{}
\begin{minipage}[b]{\linewidth}\raggedright
Term
\end{minipage} & \begin{minipage}[b]{\linewidth}\raggedright
Definition
\end{minipage} & \begin{minipage}[b]{\linewidth}\raggedright
Context
\end{minipage} \\
\midrule\noalign{}
\endhead
\bottomrule\noalign{}
\endlastfoot
\textbf{Active Inference} & A corollary of the Free Energy Principle
stating that agents act to fulfill the predictions of their generative
models. & Sec. \ref{sec:introduction} \\
\textbf{CACE} & Clarity, Appropriateness, Consistency, Evolvability ---
four-dimensional meta-standard for evaluating scientific terminology. &
Sec. \ref{sec:methodology} \\
\textbf{Generative Model} & A probabilistic model of how sensory data is
generated from latent causes. & Sec. \ref{sec:discussion} \\
\textbf{Markov Blanket} & The statistical boundary that separates
independent internal states from external states, formally defining the
individual. & Sec. \ref{sec:supplemental_analysis} \\
\textbf{Semantic Entropy} & Shannon entropy \(H(t)\) over the cluster
distribution of a term's usage contexts; quantifies terminological
ambiguity. & Sec. \ref{sec:methodology} \\
\textbf{Stigmergy} & A mechanism of indirect coordination where agents
modify the environment to stimulate the actions of others. & Sec.
\ref{sec:introduction} \\
\textbf{Superorganism} & A colony-level entity whose Markov Blanket
encompasses multiple organisms; not merely metaphorical but a formal
individuality claim. & Sec. \ref{sec:introduction}; Sec.
\ref{sec:experimental_results} \\
\textbf{Variational Free Energy} & An information-theoretic quantity
that bounds the surprise of a model; biological systems minimize this to
maintain integrity. & Sec. \ref{sec:discussion} \\
\end{longtable}
}

\subsection{Pipeline Modules}\label{pipeline-modules}

{\def\LTcaptype{none} % do not increment counter
\begin{longtable}[]{@{}
  >{\raggedright\arraybackslash}p{(\linewidth - 4\tabcolsep) * \real{0.3333}}
  >{\raggedright\arraybackslash}p{(\linewidth - 4\tabcolsep) * \real{0.3333}}
  >{\raggedright\arraybackslash}p{(\linewidth - 4\tabcolsep) * \real{0.3333}}@{}}
\toprule\noalign{}
\begin{minipage}[b]{\linewidth}\raggedright
Module
\end{minipage} & \begin{minipage}[b]{\linewidth}\raggedright
File
\end{minipage} & \begin{minipage}[b]{\linewidth}\raggedright
Function
\end{minipage} \\
\midrule\noalign{}
\endhead
\bottomrule\noalign{}
\endlastfoot
Text Processing & \texttt{src/analysis/text\_analysis.py} &
Tokenization, normalization, feature extraction \\
Term Extraction & \texttt{src/analysis/term\_extraction.py} &
Domain-aware terminology identification \\
Semantic Entropy & \texttt{src/analysis/semantic\_entropy.py} & Per-term
\(H(t)\) computation via TF-IDF + \(k\)-means \\
CACE Scoring & \texttt{src/analysis/cace\_scoring.py} & Four-dimensional
terminology evaluation \\
Domain Analysis & \texttt{src/analysis/domain\_analysis.py} & Per-domain
framing and ambiguity analysis \\
Conceptual Mapping & \texttt{src/analysis/conceptual\_mapping.py} &
Cross-domain concept graph construction \\
Rhetorical Analysis & \texttt{src/analysis/rhetorical\_analysis.py} &
Framing detection and argumentative scoring \\
Discourse Analysis & \texttt{src/analysis/discourse\_analysis.py} &
Discourse pattern classification \\
Statistics & \texttt{src/analysis/statistics.py} & Statistical
validation utilities \\
Visualization & \texttt{src/visualization/concept\_visualization.py} &
Network and domain-specific figure generation \\
\end{longtable}
}

\newpage

\section{References}\label{sec:references}

\bibliography{references}

\newpage

\section{Supplemental Methods: Text Processing and Term
Extraction}\label{sec:supplemental_methods}

This supplement documents the implementation architecture of the
Ento-Linguistic analysis pipeline. Every entry corresponds to a real
module, class, or function in \texttt{src/}. All corpus statistics cited
here are sourced from the live pipeline output in \texttt{output/data/}
and are regenerated on each clean-slate pipeline run.

\begin{center}\rule{0.5\linewidth}{0.5pt}\end{center}

\subsection{Package Architecture}\label{package-architecture}

\begin{Shaded}
\begin{Highlighting}[]
\NormalTok{src/}
\NormalTok{├── analysis/}
\NormalTok{│   ├── cace\_scoring.py         \# CACE dimension scoring}
\NormalTok{│   ├── conceptual\_mapping.py   \# Concept map construction \& analysis}
\NormalTok{│   ├── discourse\_analysis.py   \# Discourse{-}level analysis}
\NormalTok{│   ├── discourse\_patterns.py   \# Discourse pattern detection}
\NormalTok{│   ├── domain\_analysis.py      \# Six{-}domain specialist analysis}
\NormalTok{│   ├── performance.py          \# Pipeline performance metrics}
\NormalTok{│   ├── persuasive\_analysis.py  \# Persuasive strategy analysis}
\NormalTok{│   ├── rhetorical\_analysis.py  \# Rhetorical strategy \& narrative analysis}
\NormalTok{│   ├── semantic\_entropy.py     \# Semantic entropy H(t) computation}
\NormalTok{│   ├── statistics.py           \# Statistical tests (t{-}test, ANOVA, CI)}
\NormalTok{│   ├── term\_extraction.py      \# Term extraction \& classification}
\NormalTok{│   └── text\_analysis.py        \# Text normalization \& tokenization}
\NormalTok{├── core/}
\NormalTok{│   ├── exceptions.py           \# Custom exception hierarchy}
\NormalTok{│   ├── logging.py              \# Logging infrastructure}
\NormalTok{│   ├── markdown\_integration.py \# Manuscript markdown integration}
\NormalTok{│   ├── metrics.py              \# Pipeline metrics collection}
\NormalTok{│   ├── parameters.py           \# Configurable pipeline parameters}
\NormalTok{│   ├── validation.py           \# Input validation}
\NormalTok{│   └── validation\_utils.py     \# Validation helpers}
\NormalTok{├── data/}
\NormalTok{│   ├── data\_generator.py       \# Synthetic data generation for testing}
\NormalTok{│   ├── data\_processing.py      \# Data loading and transformation}
\NormalTok{│   ├── literature\_mining.py    \# Literature corpus mining}
\NormalTok{│   └── loader.py               \# Corpus file loader}
\NormalTok{├── pipeline/}
\NormalTok{│   ├── reporting.py            \# Pipeline output reporting}
\NormalTok{│   └── simulation.py           \# Simulation framework}
\NormalTok{└── visualization/}
\NormalTok{    ├── concept\_visualization.py     \# Multi{-}panel concept figures}
\NormalTok{    ├── figure\_manager.py            \# Figure registry \& integrity}
\NormalTok{    ├── plots.py                     \# Low{-}level plot utilities}
\NormalTok{    ├── statistical\_visualization.py \# Statistical plots}
\NormalTok{    └── visualization.py            \# Visualization utilities}
\end{Highlighting}
\end{Shaded}

\begin{center}\rule{0.5\linewidth}{0.5pt}\end{center}

\subsection{\texorpdfstring{Text Processing
(\texttt{src/analysis/text\_analysis.py})}{Text Processing (src/analysis/text\_analysis.py)}}\label{text-processing-srcanalysistext_analysis.py}

\subsubsection{\texorpdfstring{\texttt{TextProcessor}}{TextProcessor}}\label{textprocessor}

\textbf{Constructor parameters:}

{\def\LTcaptype{none} % do not increment counter
\begin{longtable}[]{@{}
  >{\raggedright\arraybackslash}p{(\linewidth - 6\tabcolsep) * \real{0.2821}}
  >{\raggedright\arraybackslash}p{(\linewidth - 6\tabcolsep) * \real{0.1538}}
  >{\raggedright\arraybackslash}p{(\linewidth - 6\tabcolsep) * \real{0.2308}}
  >{\raggedright\arraybackslash}p{(\linewidth - 6\tabcolsep) * \real{0.3333}}@{}}
\toprule\noalign{}
\begin{minipage}[b]{\linewidth}\raggedright
Parameter
\end{minipage} & \begin{minipage}[b]{\linewidth}\raggedright
Type
\end{minipage} & \begin{minipage}[b]{\linewidth}\raggedright
Default
\end{minipage} & \begin{minipage}[b]{\linewidth}\raggedright
Description
\end{minipage} \\
\midrule\noalign{}
\endhead
\bottomrule\noalign{}
\endlastfoot
\texttt{language} & \texttt{str} & \texttt{"english"} & NLTK processing
language \\
\texttt{custom\_stop\_words} & \texttt{Optional{[}Set{[}str{]}{]}} &
\texttt{None} & Additional domain stop-words \\
\end{longtable}
}

Stop-word vocabulary = NLTK English stop-words ∪
\texttt{SCIENTIFIC\_STOP\_WORDS} (24 domain meta-language tokens:
\emph{fig, table, et, al, etc, ie, eg, vs, cf, respectively, however,
therefore, thus, although, whereas, furthermore, moreover, addition,
similarly, consequently, subsequently, accordingly, nevertheless,
nonetheless}).

Scientific term protection vocabulary (preserved against tokenization
splitting): \emph{superorganism, eusocial, eusociality, hymenoptera,
formicidae, myrmicinae, ponerinae, dorylinae, phylogenetic, ontogenetic,
phenotypic, genotypic}.

\textbf{Methods:}

{\def\LTcaptype{none} % do not increment counter
\begin{longtable}[]{@{}
  >{\raggedright\arraybackslash}p{(\linewidth - 6\tabcolsep) * \real{0.2286}}
  >{\raggedright\arraybackslash}p{(\linewidth - 6\tabcolsep) * \real{0.3143}}
  >{\raggedright\arraybackslash}p{(\linewidth - 6\tabcolsep) * \real{0.2571}}
  >{\raggedright\arraybackslash}p{(\linewidth - 6\tabcolsep) * \real{0.2000}}@{}}
\toprule\noalign{}
\begin{minipage}[b]{\linewidth}\raggedright
Method
\end{minipage} & \begin{minipage}[b]{\linewidth}\raggedright
Signature
\end{minipage} & \begin{minipage}[b]{\linewidth}\raggedright
Returns
\end{minipage} & \begin{minipage}[b]{\linewidth}\raggedright
Notes
\end{minipage} \\
\midrule\noalign{}
\endhead
\bottomrule\noalign{}
\endlastfoot
\texttt{normalize\_text} & \texttt{(text:\ str)\ →\ str} & Normalized
string & NFKC → lowercase → punctuation removal (retaining hyphens) →
whitespace collapse \\
\texttt{tokenize\_sentences} & \texttt{(text:\ str)\ →\ List{[}str{]}} &
Sentence list & NLTK \texttt{sent\_tokenize} \\
\texttt{tokenize\_words} &
\texttt{(text:\ str,\ preserve\_scientific:\ bool)\ →\ List{[}str{]}} &
Token list & NLTK \texttt{word\_tokenize} + sliding-window
scientific-term merge \\
\texttt{remove\_punctuation} &
\texttt{(tokens:\ List{[}str{]})\ →\ List{[}str{]}} & Clean tokens &
Regex \texttt{{[}\^{}\textbackslash{}w\textbackslash{}-\_{]}} removal;
retains alphanumeric content \\
\texttt{remove\_stop\_words} &
\texttt{(tokens:\ List{[}str{]})\ →\ List{[}str{]}} & Filtered tokens &
Lowercased lookup against combined stop-word set \\
\texttt{lemmatize\_tokens} &
\texttt{(tokens:\ List{[}str{]})\ →\ List{[}str{]}} & Lemmatized tokens
& NLTK \texttt{WordNetLemmatizer.lemmatize} \\
\texttt{process\_text} &
\texttt{(text:\ str,\ lemmatize=True,\ remove\_stops=True)\ →\ List{[}str{]}}
& Processed tokens & Full pipeline: normalize → tokenize → remove punct
→ {[}stop removal{]} → {[}lemmatize{]} \\
\texttt{extract\_ngrams} &
\texttt{(tokens,\ n=2,\ min\_freq=1)\ →\ Dict{[}str,int{]}} & N-gram
counts & Sliding window; filters by \texttt{min\_freq} \\
\texttt{get\_vocabulary\_stats} &
\texttt{(texts:\ List{[}str{]})\ →\ Dict} & Stats dict & total\_tokens,
unique\_tokens, total\_characters, avg\_token\_length,
most\_common\_tokens (top 20), type\_token\_ratio \\
\end{longtable}
}

\textbf{Corpus vocabulary statistics (current run, sourced from
\texttt{output/data/corpus\_statistics.json}):}

{\def\LTcaptype{none} % do not increment counter
\begin{longtable}[]{@{}
  >{\raggedright\arraybackslash}p{(\linewidth - 2\tabcolsep) * \real{0.5333}}
  >{\raggedright\arraybackslash}p{(\linewidth - 2\tabcolsep) * \real{0.4667}}@{}}
\toprule\noalign{}
\begin{minipage}[b]{\linewidth}\raggedright
Metric
\end{minipage} & \begin{minipage}[b]{\linewidth}\raggedright
Value
\end{minipage} \\
\midrule\noalign{}
\endhead
\bottomrule\noalign{}
\endlastfoot
Total tokens & \textbf{999905} \\
Unique token types & \textbf{47227} \\
Type--token ratio & \textbf{0.0472} \\
Top 5 tokens & ant (20538), specie (10623), colony (7456), worker
(6396), social (4790) \\
\end{longtable}
}

\subsubsection{\texorpdfstring{\texttt{LinguisticFeatureExtractor}}{LinguisticFeatureExtractor}}\label{linguisticfeatureextractor}

Regex-based framing feature extraction. Three pattern sets (16 patterns
total):

\begin{itemize}
\tightlist
\item
  \textbf{Anthropomorphic} (4 patterns):
  \texttt{\textbackslash{}b(choose\textbar{}decide\textbar{}prefer\textbar{}select\textbar{}opt)\textbackslash{}b},
  \texttt{\textbackslash{}b(communicate\textbar{}signal\textbar{}inform\textbar{}warn)\textbackslash{}b},
  \texttt{\textbackslash{}b(cooperate\textbar{}compete\textbar{}negotiate\textbar{}trade)\textbackslash{}b},
  \texttt{\textbackslash{}b(recognize\textbar{}identify\textbar{}distinguish\textbar{}know)\textbackslash{}b}
\item
  \textbf{Hierarchical} (4 patterns):
  \texttt{\textbackslash{}b(superior\textbar{}inferior\textbar{}dominant\textbar{}subordinate)\textbackslash{}b},
  \texttt{\textbackslash{}b(command\textbar{}control\textbar{}authority\textbar{}obey)\textbackslash{}b},
  \texttt{\textbackslash{}b(leader\textbar{}follower\textbar{}boss\textbar{}worker)\textbackslash{}b},
  \texttt{\textbackslash{}b(ruler\textbar{}subject\textbar{}governor\textbar{}citizen)\textbackslash{}b}
\item
  \textbf{Economic} (4 patterns):
  \texttt{\textbackslash{}b(invest\textbar{}profit\textbar{}cost\textbar{}benefit)\textbackslash{}b},
  \texttt{\textbackslash{}b(trade\textbar{}exchange\textbar{}transaction\textbar{}market)\textbackslash{}b},
  \texttt{\textbackslash{}b(resource\textbar{}allocation\textbar{}distribution\textbar{}share)\textbackslash{}b},
  \texttt{\textbackslash{}b(value\textbar{}worth\textbar{}price\textbar{}commodity)\textbackslash{}b}
\end{itemize}

\texttt{extract\_framing\_features(text)} → dict with raw counts +
normalized densities (count / total\_words).

Additional methods: \texttt{detect\_terminology\_patterns(tokens)} →
compound terms, hyphenated terms, scientific abbreviations (≥2 uppercase
letters), Latin indicator tokens;
\texttt{analyze\_sentence\_complexity(text)} → sentence count, avg
sentence length, complexity ratio (sentences containing
coordinating/subordinating conjunctions or commas).

\begin{center}\rule{0.5\linewidth}{0.5pt}\end{center}

\subsection{\texorpdfstring{Terminology Extraction
(\texttt{src/analysis/term\_extraction.py})}{Terminology Extraction (src/analysis/term\_extraction.py)}}\label{terminology-extraction-srcanalysisterm_extraction.py}

\subsubsection{\texorpdfstring{\texttt{Term}
Dataclass}{Term Dataclass}}\label{term-dataclass}

\begin{Shaded}
\begin{Highlighting}[]
\AttributeTok{@dataclass}
\KeywordTok{class}\NormalTok{ Term:}
\NormalTok{    text: }\BuiltInTok{str}               \CommentTok{\# Surface form}
\NormalTok{    lemma: }\BuiltInTok{str}              \CommentTok{\# WordNet lemma}
\NormalTok{    domains: List[}\BuiltInTok{str}\NormalTok{]      }\CommentTok{\# Ento{-}Linguistic domain list}
\NormalTok{    frequency: }\BuiltInTok{int}          \CommentTok{\# Corpus{-}wide occurrence count}
\NormalTok{    contexts: List[}\BuiltInTok{str}\NormalTok{]     }\CommentTok{\# Deduplicated context sentences}
\NormalTok{    pos\_tags: List[}\BuiltInTok{str}\NormalTok{]     }\CommentTok{\# Part{-}of{-}speech tags}
\NormalTok{    confidence: }\BuiltInTok{float}       \CommentTok{\# Extraction confidence}
\NormalTok{    semantic\_entropy: }\BuiltInTok{float} \CommentTok{\# Shannon entropy H(t) in bits}
\end{Highlighting}
\end{Shaded}

Serialization: \texttt{to\_dict()} / \texttt{from\_dict()} (backward
compatible; injects \texttt{semantic\_entropy=0.0} for older records).

\subsubsection{\texorpdfstring{\texttt{TerminologyExtractor}}{TerminologyExtractor}}\label{terminologyextractor}

Domain seed lexicons (partial list):

{\def\LTcaptype{none} % do not increment counter
\begin{longtable}[]{@{}
  >{\raggedright\arraybackslash}p{(\linewidth - 2\tabcolsep) * \real{0.3478}}
  >{\raggedright\arraybackslash}p{(\linewidth - 2\tabcolsep) * \real{0.6522}}@{}}
\toprule\noalign{}
\begin{minipage}[b]{\linewidth}\raggedright
Domain
\end{minipage} & \begin{minipage}[b]{\linewidth}\raggedright
Example Seeds
\end{minipage} \\
\midrule\noalign{}
\endhead
\bottomrule\noalign{}
\endlastfoot
\texttt{unit\_of\_individuality} & ant, nestmate, colony, superorganism,
eusocial, individual, collective, organism \\
\texttt{behavior\_and\_identity} & behavior, caste, task, forager,
nurse, soldier, identity, polyethism \\
\texttt{power\_and\_labor} & queen, worker, dominance, hierarchy,
division of labor, subordinate, control \\
\texttt{sex\_and\_reproduction} & sex, reproduction, mating,
haplodiploidy, queen, egg, sperm, parthenogenesis \\
\texttt{kin\_and\_relatedness} & kin, relatedness, altruism, inclusive
fitness, nepotism, sibling \\
\texttt{economics} & cost, benefit, foraging, resource, allocation,
efficiency, trade, investment \\
\end{longtable}
}

Extraction: normalize → tokenize → match against domain seed sets →
extend via co-occurrence proximity (3-token window) → deduplicate
contexts → assign confidence.
\texttt{create\_domain\_seed\_expansion(domain\_seeds,\ corpus\_terms)}
is the domain-agnostic expansion utility.

\textbf{Pipeline run results (sourced from
\texttt{output/data/domain\_statistics.json}):}

{\def\LTcaptype{none} % do not increment counter
\begin{longtable}[]{@{}llll@{}}
\toprule\noalign{}
Domain & Term Count & Total Frequency & Bridging Terms \\
\midrule\noalign{}
\endhead
\bottomrule\noalign{}
\endlastfoot
Power \& Labor & 277 & 5621 & 139 \\
Unit of Individuality & 520 & 8095 & 24 \\
Sex \& Reproduction & 227 & 4465 & 62 \\
Behavior \& Identity & 243 & 9239 & 91 \\
Kin \& Relatedness & 72 & 2569 & 4 \\
Economics & 65 & 2009 & 3 \\
\textbf{Total (all domains)} & \textbf{1239} & --- & --- \\
\end{longtable}
}

\begin{center}\rule{0.5\linewidth}{0.5pt}\end{center}

\subsection{\texorpdfstring{Semantic Entropy
(\texttt{src/analysis/semantic\_entropy.py})}{Semantic Entropy (src/analysis/semantic\_entropy.py)}}\label{semantic-entropy-srcanalysissemantic_entropy.py}

\subsubsection{Constants}\label{constants}

\begin{Shaded}
\begin{Highlighting}[]
\NormalTok{HIGH\_ENTROPY\_THRESHOLD }\OperatorTok{=} \FloatTok{2.0}  \CommentTok{\# bits; corresponds to ≥4 equiprobable senses}
\end{Highlighting}
\end{Shaded}

\subsubsection{\texorpdfstring{\texttt{SemanticEntropyResult}
Dataclass}{SemanticEntropyResult Dataclass}}\label{semanticentropyresult-dataclass}

\begin{Shaded}
\begin{Highlighting}[]
\AttributeTok{@dataclass}
\KeywordTok{class}\NormalTok{ SemanticEntropyResult:}
\NormalTok{    term: }\BuiltInTok{str}
\NormalTok{    entropy\_bits: }\BuiltInTok{float}            \CommentTok{\# Shannon H(t) in bits (base 2)}
\NormalTok{    n\_clusters: }\BuiltInTok{int}                \CommentTok{\# KMeans k actually used}
\NormalTok{    cluster\_distribution: List[}\BuiltInTok{float}\NormalTok{]  }\CommentTok{\# Empirical p(c\_i) per cluster}
\NormalTok{    is\_high\_entropy: }\BuiltInTok{bool}          \CommentTok{\# True if entropy\_bits \textgreater{} 2.0}
\NormalTok{    n\_contexts: }\BuiltInTok{int}                \CommentTok{\# Valid contexts used}
\NormalTok{    h\_max: }\BuiltInTok{float}                   \CommentTok{\# log2(k): maximum attainable entropy for k clusters}
\NormalTok{    entropy\_normalized: }\BuiltInTok{float}      \CommentTok{\# H / h\_max, normalized to [0, 1]}
\end{Highlighting}
\end{Shaded}

\subsubsection{\texorpdfstring{\texttt{calculate\_semantic\_entropy}}{calculate\_semantic\_entropy}}\label{calculate_semantic_entropy}

\begin{Shaded}
\begin{Highlighting}[]
\KeywordTok{def}\NormalTok{ calculate\_semantic\_entropy(}
\NormalTok{    term: }\BuiltInTok{str}\NormalTok{,}
\NormalTok{    contexts: List[}\BuiltInTok{str}\NormalTok{],}
\NormalTok{    max\_clusters: }\BuiltInTok{int} \OperatorTok{=} \DecValTok{5}\NormalTok{,}
\NormalTok{    min\_contexts: }\BuiltInTok{int} \OperatorTok{=} \DecValTok{5}\NormalTok{,}
\NormalTok{    random\_state: }\BuiltInTok{int} \OperatorTok{=} \DecValTok{42}\NormalTok{,}
\NormalTok{    threshold: }\BuiltInTok{float} \OperatorTok{=} \FloatTok{2.0}\NormalTok{,}
\NormalTok{) }\OperatorTok{{-}\textgreater{}}\NormalTok{ SemanticEntropyResult}
\end{Highlighting}
\end{Shaded}

\textbf{Algorithm:}

\begin{enumerate}
\def\labelenumi{\arabic{enumi}.}
\tightlist
\item
  Filter to contexts with ≥3 whitespace-delimited words.
\item
  If valid contexts \textless{} \texttt{min\_contexts}: return H=0.0,
  n\_clusters=1 (or 0 if empty).
\item
  TF-IDF:
  \texttt{TfidfVectorizer(stop\_words="english",\ min\_df=1,\ max\_features=1000)}.
\item
  Cluster count:
  \texttt{k\_upper\ =\ min(max\_clusters,\ n\ -\ 1,\ max(2,\ int(np.sqrt(n))))};
  \texttt{k\ =\ max(2,\ k\_upper)} with
  \texttt{n\ =\ len(valid\_contexts)} --- guarantees
  \texttt{k\ \textless{}\ n} and caps \texttt{k} at the data-driven
  \texttt{sqrt(n)} bound; if \texttt{k\ \textless{}\ 2} return H=0.0.
\item
  \texttt{KMeans(n\_clusters=k,\ random\_state=42,\ n\_init=10)} →
  labels.
\item
  Empirical distribution: \(p_i = n_i / N\).
\item
  \(H =\) \texttt{scipy.stats.entropy(probabilities,\ base=2)}.
\item
  Normalization: \texttt{h\_max\ =\ log2(k)};
  \texttt{entropy\_normalized\ =\ entropy\_bits\ /\ h\_max} (0.0 when
  \texttt{k\ \textless{}\ 2}).
\item
  Exception guard: any sklearn/scipy failure → H=0.0.
\end{enumerate}

\subsubsection{Corpus-Level Functions}\label{corpus-level-functions}

{\def\LTcaptype{none} % do not increment counter
\begin{longtable}[]{@{}
  >{\raggedright\arraybackslash}p{(\linewidth - 4\tabcolsep) * \real{0.3125}}
  >{\raggedright\arraybackslash}p{(\linewidth - 4\tabcolsep) * \real{0.2812}}
  >{\raggedright\arraybackslash}p{(\linewidth - 4\tabcolsep) * \real{0.4062}}@{}}
\toprule\noalign{}
\begin{minipage}[b]{\linewidth}\raggedright
Function
\end{minipage} & \begin{minipage}[b]{\linewidth}\raggedright
Returns
\end{minipage} & \begin{minipage}[b]{\linewidth}\raggedright
Description
\end{minipage} \\
\midrule\noalign{}
\endhead
\bottomrule\noalign{}
\endlastfoot
\texttt{calculate\_corpus\_entropy(terms\_contexts,\ ...)} &
\texttt{Dict{[}str,\ SemanticEntropyResult{]}} & Runs per-term entropy
for all terms \\
\texttt{get\_high\_entropy\_terms(results)} &
\texttt{List{[}SemanticEntropyResult{]}} & Filters
\texttt{is\_high\_entropy=True}, sorted descending by
\texttt{entropy\_bits} \\
\end{longtable}
}

\newpage

\section{Supplemental Methods: Statistical and Scoring
Infrastructure}\label{sec:supplemental_infrastructure}

\subsection{\texorpdfstring{Statistical Analysis
(\texttt{src/analysis/statistics.py})}{Statistical Analysis (src/analysis/statistics.py)}}\label{statistical-analysis-srcanalysisstatistics.py}

All functions implemented from mathematical first principles via NumPy
and SciPy.

\subsubsection{\texorpdfstring{\texttt{DescriptiveStats}}{DescriptiveStats}}\label{descriptivestats}

\texttt{calculate\_descriptive\_stats(data:\ np.ndarray)} →
\texttt{DescriptiveStats(mean,\ std,\ median,\ min,\ max,\ q25,\ q75,\ count)}.

\subsubsection{\texorpdfstring{\texttt{t\_test}}{t\_test}}\label{t_test}

\begin{Shaded}
\begin{Highlighting}[]
\KeywordTok{def}\NormalTok{ t\_test(sample1, sample2}\OperatorTok{=}\VariableTok{None}\NormalTok{, mu}\OperatorTok{=}\VariableTok{None}\NormalTok{, alternative}\OperatorTok{=}\StringTok{"two{-}sided"}\NormalTok{) }\OperatorTok{{-}\textgreater{}}\NormalTok{ Dict}
\CommentTok{\# Returns: t\_statistic, p\_value, degrees\_of\_freedom, alternative}
\end{Highlighting}
\end{Shaded}

\begin{itemize}
\tightlist
\item
  One-sample: \(t = (\bar{x} - \mu_0) / (s / \sqrt{n})\), \(df = n-1\)
\item
  Two-sample (Welch):
  \(t = (\bar{x}_1 - \bar{x}_2) / \sqrt{s_1^2/n_1 + s_2^2/n_2}\);
  Welch--Satterthwaite df
\item
  \(p\)-value via \texttt{scipy.stats.t.sf}
\end{itemize}

\subsubsection{\texorpdfstring{\texttt{anova\_test}}{anova\_test}}\label{anova_test}

\texttt{anova\_test(groups:\ List{[}np.ndarray{]})} →
\texttt{f\_statistic,\ p\_value,\ df\_between,\ df\_within}.
From-scratch SS computation; \(p\) via \texttt{scipy.stats.f.sf}.

\subsubsection{\texorpdfstring{\texttt{calculate\_correlation}}{calculate\_correlation}}\label{calculate_correlation}

\texttt{method="pearson"}: \texttt{numpy.corrcoef} + t-distribution
p-value. \texttt{method="spearman"}: \texttt{scipy.stats.spearmanr}.

\subsubsection{\texorpdfstring{\texttt{calculate\_confidence\_interval}}{calculate\_confidence\_interval}}\label{calculate_confidence_interval}

\(\bar{x} \pm t_{0.975, n-1} \cdot s/\sqrt{n}\); critical value via
\texttt{scipy.stats.t.ppf(0.975,\ n-1)}.

\subsubsection{\texorpdfstring{\texttt{fit\_distribution}}{fit\_distribution}}\label{fit_distribution}

Supports \texttt{"normal"} (MLE: μ, σ), \texttt{"exponential"} (MLE:
λ=1/mean), \texttt{"uniform"} (min, max).

\begin{center}\rule{0.5\linewidth}{0.5pt}\end{center}

\subsection{\texorpdfstring{Domain Analysis
(\texttt{src/analysis/domain\_analysis.py})}{Domain Analysis (src/analysis/domain\_analysis.py)}}\label{domain-analysis-srcanalysisdomain_analysis.py}

\subsubsection{\texorpdfstring{\texttt{DomainAnalysis}
Dataclass}{DomainAnalysis Dataclass}}\label{domainanalysis-dataclass}

Fields: \texttt{domain\_name}, \texttt{key\_terms} (top-10 by
frequency), \texttt{term\_patterns}
(compound/multi\_word/capitalized/abbreviation/numeric counts),
\texttt{framing\_assumptions}, \texttt{conceptual\_structure} (domain
ontology), \texttt{ambiguities} (term/contexts/issue triplets),
\texttt{recommendations}, \texttt{frequency\_stats}
(mean/median/SD/histogram), \texttt{cooccurrence\_analysis},
\texttt{ambiguity\_metrics}, \texttt{confidence\_scores},
\texttt{conceptual\_metrics}, \texttt{statistical\_significance}.

\subsubsection{\texorpdfstring{\texttt{DomainAnalyzer}
Methods}{DomainAnalyzer Methods}}\label{domainanalyzer-methods}

\texttt{analyze\_all\_domains(terms,\ texts)} dispatches to six
specialist methods per domain, then enriches with:

\begin{enumerate}
\def\labelenumi{\arabic{enumi}.}
\tightlist
\item
  \texttt{analyze\_term\_frequency\_distribution}: NumPy
  \texttt{histogram(bins="auto")}; top-10 term--frequency pairs.
\item
  \texttt{analyze\_term\_cooccurrence}: sliding-window co-occurrence
  matrix.
\item
  \texttt{quantify\_ambiguity\_metrics}: domain-level semantic entropy
  aggregation.
\item
  \texttt{calculate\_statistical\_significance}: \(\chi^2\)/Fisher's on
  pattern distributions.
\end{enumerate}

Term pattern counting (\texttt{\_analyze\_term\_patterns}): compound
(contains \texttt{\_}/\texttt{-}), multi\_word (contains \texttt{}),
capitalized, abbreviation (\texttt{\^{}{[}A-Z{]}\{2,\}\$}), numeric.

\begin{center}\rule{0.5\linewidth}{0.5pt}\end{center}

\subsection{\texorpdfstring{Conceptual Mapping
(\texttt{src/analysis/conceptual\_mapping.py})}{Conceptual Mapping (src/analysis/conceptual\_mapping.py)}}\label{conceptual-mapping-srcanalysisconceptual_mapping.py}

\subsubsection{Data Structures}\label{data-structures}

\texttt{Concept}: name, description, terms (Set), domains (Set),
parent\_concepts, child\_concepts, confidence. \texttt{ConceptMap}:
\texttt{concepts:\ Dict{[}str,\ Concept{]}},
\texttt{term\_to\_concepts:\ Dict{[}str,\ Set{[}str{]}{]}},
\texttt{concept\_relationships:\ Dict{[}Tuple{[}str,str{]},\ float{]}}.

\subsubsection{\texorpdfstring{\texttt{ConceptualMapper}}{ConceptualMapper}}\label{conceptualmapper}

\texttt{build\_concept\_map(terms)}: (1) instantiate 6 base concept
nodes; (2) domain- and keyword-based term assignment; (3)
overlap-coefficient edge creation.

\texttt{analyze\_concept\_centrality}: NetworkX
degree/betweenness/closeness/eigenvector centrality (pure-Python
fallback). \texttt{quantify\_relationship\_strength}: composite =
base×0.4 + term\_overlap×0.3 + domain\_overlap×0.2 + hierarchical×0.1.
\texttt{identify\_cross\_domain\_bridges}: concepts spanning ≥2 domains.
\texttt{calculate\_concept\_similarity}: Jaccard + domain overlap bonus
(max 0.3). \texttt{detect\_anthropomorphic\_concepts}: 5 indicator
categories (agency/communication/social\_contract/cognition/hierarchy).

\textbf{Pipeline results (sourced from
\texttt{output/data/concept\_map\_summary.json}):}

{\def\LTcaptype{none} % do not increment counter
\begin{longtable}[]{@{}lll@{}}
\toprule\noalign{}
Concept & Terms & Domains \\
\midrule\noalign{}
\endhead
\bottomrule\noalign{}
\endlastfoot
\texttt{biological\_individuality} & 568 & Unit of Individuality \\
\texttt{social\_organization} & 556 & Power \& Labor; Behavior \&
Identity \\
\texttt{reproductive\_biology} & 256 & Sex \& Reproduction \\
\texttt{kinship\_systems} & 311 & Kin \& Relatedness \\
\texttt{resource\_economics} & 96 & Economics \\
\texttt{behavioral\_ecology} & 376 & Behavior \& Identity; Economics \\
\textbf{Concept relationships} & \textbf{15} & \\
\end{longtable}
}

\begin{center}\rule{0.5\linewidth}{0.5pt}\end{center}

\subsection{\texorpdfstring{CACE Scoring
(\texttt{src/analysis/cace\_scoring.py})}{CACE Scoring (src/analysis/cace\_scoring.py)}}\label{cace-scoring-srcanalysiscace_scoring.py}

\texttt{CACEScore}: term, clarity, appropriateness, consistency,
evolvability, aggregate (mean of four).

{\def\LTcaptype{none} % do not increment counter
\begin{longtable}[]{@{}
  >{\raggedright\arraybackslash}p{(\linewidth - 2\tabcolsep) * \real{0.5263}}
  >{\raggedright\arraybackslash}p{(\linewidth - 2\tabcolsep) * \real{0.4737}}@{}}
\toprule\noalign{}
\begin{minipage}[b]{\linewidth}\raggedright
Function
\end{minipage} & \begin{minipage}[b]{\linewidth}\raggedright
Formula
\end{minipage} \\
\midrule\noalign{}
\endhead
\bottomrule\noalign{}
\endlastfoot
\texttt{score\_clarity} &
\texttt{max(0,\ 1\ -\ entropy\_bits\ /\ log2(10))} ---
\texttt{log2(10)\ ≈\ 3.32} = \texttt{DEFAULT\_MAX\_ENTROPY} \\
\texttt{score\_appropriateness} &
\texttt{1\ -\ (0.4\ ×\ 𝟙{[}term\ ∩\ 𝒜{]}\ +\ 0.1\ ×\ overlap\ +\ 0.05\ ×\ max(domains−1,\ 0))}
--- graduated penalty, never zeroed \\
\texttt{score\_consistency} & Mean pairwise cosine similarity of TF-IDF
context vectors (high = consistent) \\
\texttt{score\_evolvability} &
\texttt{0.5\ ×\ min(1,\ domains/3)\ +\ 0.5\ ×\ min(1,\ scale\_levels\_in\_contexts/3)} \\
\texttt{evaluate\_term\_cace} & All four scorers → \texttt{CACEScore} \\
\texttt{compare\_terms\_cace} & Ranked \texttt{List{[}CACEScore{]}} by
aggregate descending \\
\end{longtable}
}

\texttt{ANTHROPOMORPHIC\_TERMS}: queen, king, slave, worker, soldier,
nurse, princess, maiden, + additional (full set in source).

\begin{center}\rule{0.5\linewidth}{0.5pt}\end{center}

\subsection{\texorpdfstring{Rhetorical Analysis
(\texttt{src/analysis/rhetorical\_analysis.py})}{Rhetorical Analysis (src/analysis/rhetorical\_analysis.py)}}\label{rhetorical-analysis-srcanalysisrhetorical_analysis.py}

\texttt{analyze\_rhetorical\_strategies}: 4 strategy types,
regex-detected per abstract (authority, analogy, generalization,
anecdotal). \texttt{identify\_narrative\_frameworks}: 4 framework types
(progress/conflict/discovery/complexity), keyword-presence classifier.
\texttt{quantify\_rhetorical\_patterns}: total\_occurrences,
text\_coverage, effectiveness = min(occurrences/n\_texts, 1),
persuasiveness. \texttt{score\_argumentative\_structures}:
claim\_strength + evidence\_quality + reasoning\_coherence (mean) +
confidence\_score. \texttt{analyze\_narrative\_frequency}: frequency,
coverage\_percentage, avg\_text\_length, unique\_bigram\_count,
consistency\_score.

\begin{center}\rule{0.5\linewidth}{0.5pt}\end{center}

\subsection{\texorpdfstring{Visualization
(\texttt{src/visualization/})}{Visualization (src/visualization/)}}\label{visualization-srcvisualization}

\texttt{ConceptVisualizer} generates 11 research figures via matplotlib
multi-panel layouts. \texttt{FigureManager} maintains a JSON figure
registry with SHA integrity hashes. \texttt{StatisticalVisualization}
produces forest plots, violin plots, heatmaps, and regression
diagnostics.

\textbf{Current run:} figures are generated by the visualization
pipeline; \texttt{FigureManager} records a registry entry and integrity
hash per deliverable when the pipeline completes successfully.

\begin{center}\rule{0.5\linewidth}{0.5pt}\end{center}

\subsection{\texorpdfstring{Core Infrastructure
(\texttt{src/core/})}{Core Infrastructure (src/core/)}}\label{core-infrastructure-srccore}

\texttt{parameters.py}: \texttt{ParameterSet}/\texttt{ParameterSweep}
infrastructure plus \texttt{AnalysisParameters} defaults
(\texttt{min\_term\_freq=5}, \texttt{cooccurrence\_threshold=2}). The
semantic-entropy parameters (\texttt{max\_clusters=5},
\texttt{min\_contexts=5}, \texttt{threshold=2.0},
\texttt{random\_state=42}) are function defaults of
\texttt{calculate\_semantic\_entropy}; TF-IDF
\texttt{max\_features=1000} is set in the vectorizer. Context windows
are per-call arguments: 3-token context-extraction windows in
\texttt{term\_extraction.py}, 10-word term co-occurrence windows
(\texttt{TerminologyExtractor.find\_term\_cooccurrences}), and 50-word
domain co-occurrence windows
(\texttt{DomainAnalyzer.analyze\_term\_cooccurrence}).
\texttt{validation.py}/\texttt{validation\_utils.py}: type checks and
domain membership guards on all public API entry points.
\texttt{metrics.py}: wall-clock, memory, throughput per stage.
\texttt{markdown\_integration.py}: \texttt{\textbackslash{}ref\{\}}
resolution and cross-reference validation.

\begin{center}\rule{0.5\linewidth}{0.5pt}\end{center}

\subsection{Reproducibility}\label{reproducibility}

\begin{itemize}
\tightlist
\item
  \textbf{Deterministic}: \texttt{random\_state=42} in all KMeans calls.
\item
  \textbf{Clean-slate}: \texttt{output/figures/} and
  \texttt{output/data/} wiped and recreated on every run
  (\texttt{\_setup\_directories} in
  \texttt{scripts/02\_generate\_figures.py}).
\item
  \textbf{Live statistics}: all corpus metrics are read from
  \texttt{output/data/corpus\_statistics.json},
  \texttt{domain\_statistics.json}, and
  \texttt{concept\_map\_summary.json} and substituted into the
  manuscript prose at PDF build time via template placeholders
  (\texttt{scripts/\_render\_pdf\_override.py}) --- no corpus statistic
  appears as a hardcoded literal in the manuscript.
\item
  \textbf{Dependency pinning}: all Python dependencies pinned in
  \texttt{pyproject.toml}.
\item
  \textbf{Test suite}: comprehensive test suite covering all
  \texttt{src/} modules; run via
  \texttt{uv\ run\ pytest\ tests/\ -\/-cov=src} from the project root.
\end{itemize}

\newpage

\section{Supplemental Results}\label{sec:supplemental_results}

\subsection{Pairwise Domain
Comparisons}\label{pairwise-domain-comparisons}

Table \ref{tab:pairwise_domain} presents pairwise comparisons of
per-term semantic entropy between all Ento-Linguistic domains using
Welch's two-sample \(t\)-tests. Raw \(p\)-values are computed from the
\(t\)-distribution with Satterthwaite-approximated degrees of freedom;
adjusted \(p\)-values correct for 15 simultaneous comparisons using the
Benjamini-Hochberg (BH) procedure at \(q = 0.05\). Cohen's \(d\)
quantifies effect size, interpreted as small (\(d \approx 0.2\)), medium
(\(d \approx 0.5\)), or large (\(d \geq 0.8\)). Domain descriptives
entering these tests are per-term valid-entropy means with exclusions
counted.

\begin{table}[h]
\centering
\small
\begin{tabular}{|l|l|c|c|c|c|c|}
\hline
\textbf{Domain A} & \textbf{Domain B} & \textbf{$t$} & \textbf{$p$ (raw)} & \textbf{$p$ (BH)} & \textbf{Cohen's $d$} & \textbf{Sig.\ (BH)} \\
\hline
Behavior \& Identity & Economics & -0.4516 & 0.6565 & 0.9232 & -0.1427 & no \\
Behavior \& Identity & Kin \& Relatedness & -0.2270 & 0.8227 & 0.9232 & -0.0726 & no \\
Behavior \& Identity & Power \& Labor & 0.2379 & 0.8126 & 0.9232 & 0.0495 & no \\
Behavior \& Identity & Sex \& Reproduction & -0.2542 & 0.8000 & 0.9232 & -0.0553 & no \\
Behavior \& Identity & Unit of Individuality & 0.6740 & 0.5028 & 0.9232 & 0.1300 & no \\
Economics & Kin \& Relatedness & 0.1763 & 0.8616 & 0.9232 & 0.0680 & no \\
Economics & Power \& Labor & 0.6437 & 0.5288 & 0.9232 & 0.1947 & no \\
Economics & Sex \& Reproduction & 0.2728 & 0.7879 & 0.9232 & 0.0810 & no \\
Economics & Unit of Individuality & 0.9452 & 0.3606 & 0.9232 & 0.2769 & no \\
Kin \& Relatedness & Power \& Labor & 0.3994 & 0.6945 & 0.9232 & 0.1234 & no \\
Kin \& Relatedness & Sex \& Reproduction & 0.0508 & 0.9600 & 0.9600 & 0.0154 & no \\
Kin \& Relatedness & Unit of Individuality & 0.6800 & 0.5069 & 0.9232 & 0.2050 & no \\
Power \& Labor & Sex \& Reproduction & -0.5315 & 0.5963 & 0.9232 & -0.1044 & no \\
Power \& Labor & Unit of Individuality & 0.4954 & 0.6212 & 0.9232 & 0.0804 & no \\
Sex \& Reproduction & Unit of Individuality & 1.0127 & 0.3142 & 0.9232 & 0.1845 & no \\
\hline
\end{tabular}
\caption{Pairwise Welch's $t$-tests on per-term semantic entropy between Ento-Linguistic domains, with Benjamini-Hochberg adjusted $p$-values (benjamini_hochberg correction over 15 comparisons) and Cohen's $d$ effect sizes. The Sig.\ (BH) column reports BH-adjusted significance at $q = 0.05$ for each comparison. Domain descriptives feeding these tests are per-term valid-entropy means with exclusions counted. The omnibus one-way ANOVA across the six domains on semantic_entropy yields $F(5.0000, 264.0000) = 0.3771$, 0.8642, $\eta^2 = 0.0071$, where $df_1 = k - 1$ (between-group) and $df_2 = N - k$ (within-group).}
\label{tab:pairwise_domain}
\end{table}

\subsection{CACE Scoring for Key
Terms}\label{cace-scoring-for-key-terms}

Table \ref{tab:cace_full} presents full CACE evaluations for a
representative set of entomological terms, comparing anthropomorphic
labels with proposed functional alternatives.

\begin{table}[h]
\centering
\small
\begin{tabular}{|l|c|c|c|c|c|}
\hline
\textbf{Term} & \textbf{Clarity} & \textbf{Appropriateness} & \textbf{Consistency} & \textbf{Evolvability} & \textbf{Aggregate} \\
\hline
queen & 0.37 & 0.45 & 0.04 & 0.33 & 0.30 \\
\textit{primary reproductive} & 1.00 & 1.00 & 0.50 & 0.00 & 0.62 \\
\hline
worker & 0.39 & 0.45 & 0.05 & 0.33 & 0.31 \\
\textit{non-reproductive helper} & 1.00 & 1.00 & 0.50 & 0.00 & 0.62 \\
\hline
slave & 1.00 & 0.50 & 0.50 & 0.00 & 0.50 \\
\textit{host worker} & 1.00 & 0.40 & 0.50 & 0.00 & 0.47 \\
\hline
caste & 0.32 & 1.00 & 0.06 & 0.50 & 0.47 \\
\textit{task group} & 1.00 & 1.00 & 0.50 & 0.00 & 0.62 \\
\hline
soldier & 0.37 & 0.50 & 0.10 & 0.17 & 0.28 \\
\textit{major worker} & 1.00 & 0.50 & 0.50 & 0.00 & 0.50 \\
\hline
colony & 0.39 & 1.00 & 0.04 & 0.33 & 0.44 \\
haplodiploidy & 0.38 & 1.00 & 0.05 & 0.17 & 0.40 \\
trophallaxis & 1.00 & 1.00 & 0.50 & 0.00 & 0.62 \\
\hline
\end{tabular}
\caption{CACE dimension scores for representative entomological terms. Anthropomorphic terms (queen, worker, slave, caste, soldier) consistently score lower than functional alternatives (italicized). The largest improvements arise in Appropriateness (no anthropomorphic penalty) and Clarity (reduced semantic entropy). Non-anthropomorphic technical terms (haplodiploidy, trophallaxis) score highest on Clarity due to unambiguous, single-sense usage. Note: ``colony'' receives Appropriateness $= 1.00$ because it falls outside the \texttt{ANTHROPOMORPHIC\_TERMS} set used for automated scoring; its colonial and settler-historical connotations are analyzed qualitatively in Section~\ref{sec:discussion}. All values resolve at render time from the frozen statistics artifact (\texttt{output/data/statistical\_analysis.json}).}
\label{tab:cace_full}
\end{table}

\subsection{Semantic Entropy
Distribution}\label{semantic-entropy-distribution}

Table \ref{tab:entropy_distribution} summarizes the distribution of
semantic entropy across domains.

\begin{table}[h]
\centering
\begin{tabular}{|l|c|c|c|}
\hline
\textbf{Domain} & \textbf{Mean $H$ (bits)} & \textbf{High-entropy terms (\%)} & \textbf{$N$} \\
\hline
Economics & 0.28 & 4.6 & 65 \\
Power \& Labor & 0.30 & 3.2 & 277 \\
Behavior \& Identity & 0.21 & 3.3 & 243 \\
Sex \& Reproduction & 0.27 & 4.4 & 227 \\
Unit of Individuality & 0.28 & 3.8 & 520 \\
Kin \& Relatedness & 0.25 & 2.8 & 72 \\
\hline
\textbf{Overall} & 0.27 & 3.7 & \textbf{ 1404 } \\
\hline
\end{tabular}
\caption{Distribution of semantic entropy $H(t)$ across Ento-Linguistic domains, computed from pipeline output in \texttt{output/data/domain\_statistics.json}. High-entropy terms are those exceeding the $H > 2.0$ bits threshold (per \texttt{src/analysis/semantic\_entropy.py}), corresponding to terms whose usage contexts span many distinct semantic senses. Entropy is calculated via TF-IDF vectorization of each term's corpus contexts followed by KMeans clustering (with $k < n$ contexts; see Eq.~\ref{eq:semantic_entropy}). The number of clusters is set to $k = \max(2,\, \min(k_{\max},\, n{-}1,\, \max(2, \lfloor\!\sqrt{n}\rfloor)))$ with $k_{\max}=5$ to prevent degenerate one-point-per-cluster assignments; each result also reports $H_{\max} = \log_2 k$ and the normalized entropy $H/H_{\max} \in [0,1]$. The $N$ column is each domain's term count in the same artifact, and the Overall row is their exact sum; every table value resolves at render time from the artifact.}
\label{tab:entropy_distribution}
\end{table}

\subsection{Confidence Intervals for Domain
Metrics}\label{confidence-intervals-for-domain-metrics}

The frozen statistics artifact
(\texttt{output/data/statistical\_analysis.json}) reports per-term
valid-entropy descriptives (means and standard deviations, with
exclusions counted) and the inferential results in Table
\ref{tab:pairwise_domain}; it does not compute domain-level confidence
intervals, so per-domain ambiguity-score and context-variability
intervals are not tabulated here. Separation between domains is
established inferentially by the Welch \(t\)-tests and the omnibus ANOVA
reported in Table \ref{tab:pairwise_domain}, with per-domain entropy
descriptives in Table \ref{tab:entropy_distribution} and the
accompanying summary figure \texttt{statistical\_analysis.png}.

\subsection{Full-Text Parallel Layer}\label{full-text-parallel-layer}

To test whether the ento-linguistic patterns reported above survive
beyond abstracts, a parallel analysis layer was run over 7070 open
access full texts harvested from PubMed Central (PMC). The abstract
corpus remains the headline corpus of this study; the full-text layer is
reported here as a robustness check. The analysis machinery is shared:
the same terminology extraction, domain assignment, semantic-entropy,
and CACE scoring implementations are applied to full texts, with the
frozen artifact written to \texttt{output/data/fulltext\_analysis.json}
and rendered in figure \texttt{fulltext\_analysis.png}.

The layer comprises 44484942 tokens of running text, with a per-document
median of 5621.0000 tokens. Table \ref{tab:fulltext_domain} reports
per-domain term counts and mean semantic entropy over the full texts; 15
pairwise Welch \(t\)-tests (Benjamini-Hochberg corrected, as in Table
\ref{tab:pairwise_domain}) accompany the omnibus one-way ANOVA on
per-term semantic entropy, \(F = 4.3251\), 0.0007.

Anthropomorphic framing over the same full texts is scored as the
proportion of domain-term occurrence contexts containing an
anthropomorphic framing marker (\texttt{add\_framing\_analysis} in
\texttt{src/pipeline/fulltext\_pipeline.py}): 0.0170 overall, with
Economics a clear outlier at 0.0957 against 0.0072 for Unit of
Individuality and 0.0064 for Behavior \& Identity.

\begin{table}[h]
\centering
\begin{tabular}{|l|c|c|}
\hline
\textbf{Domain} & \textbf{Terms extracted} & \textbf{Mean $H$ (bits)} \\
\hline
Behavior \& Identity & 77 & 2.0202 \\
Economics & 21 & 1.8650 \\
Kin \& Relatedness & 12 & 2.0317 \\
Power \& Labor & 100 & 2.0589 \\
Sex \& Reproduction & 68 & 2.0810 \\
Unit of Individuality & 230 & 2.0404 \\
\hline
\end{tabular}
\caption{Per-domain terminology in the PMC full-text parallel layer: extracted-term counts and mean semantic entropy $H(t)$, computed with the same pipeline as the abstract layer (\texttt{output/data/fulltext\_analysis.json}, \texttt{descriptives} section). Term extraction uses a higher minimum token frequency than the abstract layer because full texts are substantially longer.}
\label{tab:fulltext_domain}
\end{table}

\subsection{Discourse and Rhetorical
Layer}\label{discourse-and-rhetorical-layer}

Both statistical artifacts additionally carry a corpus-level
\texttt{discourse} section --- discourse patterns, rhetorical
strategies, argumentative structures, and persuasive techniques ---
computed by the shared discourse stage
(\texttt{add\_discourse\_analysis} in
\texttt{src/pipeline/statistics\_pipeline.py}) after the statistics
stages. The abstract layer's discourse pass covers 7592 texts (sample
fraction 1.0000). The full-text layer's discourse pass covers 1414 of
its 7070 texts --- a deterministic 0.2000 sample disclosed here because
full texts exceed the discourse pass's minimum-length bound far less
often than abstracts but are subsampled to keep the pass bounded. Both
layers' discourse dimensions are compared in figure
\texttt{discourse\_comparison.png} (registered as
\texttt{fig:discourse\_comparison}); panels spanning orders of magnitude
use a symlog frequency axis.

Discourse patterns occur 1278 times as hierarchical framing in the
abstract layer against 1108 occurrences in the full-text layer, with
1046 versus 1090 economic metaphors, 32 versus 47 anthropomorphic
framings, and 27 versus 33 scale-ambiguous constructions.

Rhetorical strategies show the same abstract-to-full-text expansion:
1580 anecdotal markers in the abstract layer versus 9229 in full texts,
247 versus 29178 authority markers, 455 versus 1792 analogies, and 1148
versus 4228 generalizations.

The argumentative-structure pass identifies 1804 argumentative
structures in the abstract layer versus 1299 in the full-text sample.
Metaphorical language, the dominant persuasive technique, occurs 3835
times in the abstract layer against 20759 occurrences in full texts. All
frequencies are raw corpus counts from the artifacts' \texttt{discourse}
sections (\texttt{output/data/statistical\_analysis.json} and
\texttt{output/data/fulltext\_analysis.json}); no values in this
subsection are literals.

\newpage

\section{Supplemental Analysis: Theoretical
Extensions}\label{sec:supplemental_analysis}

This section provides analytical results and theoretical extensions that
complement the main findings presented in Sections \ref{sec:methodology}
and \ref{sec:experimental_results}.

\subsection{Theoretical Extensions}\label{theoretical-extensions}

\subsubsection{Formalism of Individuality: Markov
Blankets}\label{formalism-of-individuality-markov-blankets}

To rigorize the ``Unit of Individuality'' domain, we employ the
\textbf{Markov Blanket} formalism
\cite{friston2013life, kirchhoff2018markov}. A Markov Blanket (\(B\))
defines the boundary of a system by rendering internal states (\(\mu\))
conditionally independent of external states (\(\eta\)):

\begin{equation}\label{eq:markov_blanket}
P(\mu | \eta, B) = P(\mu | B)
\end{equation}

In biological systems, the blanket consists of \textbf{sensory states}
(inputs) and \textbf{active states} (outputs).

\begin{itemize}
\tightlist
\item
  \textbf{Organismal Blanket}: The ant's cuticle and sensory receptors.
\item
  \textbf{Colonial Blanket}: The nest entrance, shared pheromone fields,
  and cuticular hydrocarbon profiles.
\end{itemize}

Linguistic confusion arises when terms index the wrong blanket.
``Superorganism'' is not a metaphor but a formal claim that the relevant
Markov Blanket enclosing the \textbf{generative model} is at the colony
level. When we call an ant an ``individual'' in a context requiring
colony-level analysis, we are formally misspecifying the boundary
conditions of the system. The Active Inferants framework
\cite{friedman2021active} operationalizes this insight, showing that
foraging behavior emerges from ensemble-level inference over pheromone
gradients---locating the generative model at the colony blanket rather
than the organismal blanket.

\subsubsection{Discourse Analysis
Frameworks}\label{discourse-analysis-frameworks}

Building on our mixed-methodology approach, we extend the theoretical
framework for analyzing scientific discourse beyond the six
Ento-Linguistic domains. Our analysis reveals that terminology networks
serve as both descriptive tools and constitutive elements of scientific
knowledge production.

\textbf{Constitutive Framework}:

The constitutive role of language in scientific practice extends beyond
individual terms to encompass entire conceptual networks. We formalize
this through the concept of \textbf{discursive framing networks}:

\begin{equation}\label{eq:discursive_framing}
F(D, T) = \sum_{t \in T} w_t \cdot f_t(D) \cdot c_t
\end{equation}

where \(D\) represents a domain, \(T\) is the terminology set, \(w_t\)
are term weights, \(f_t(D)\) is the framing function for domain \(D\),
and \(c_t\) represents contextual factors.

\subsubsection{Ambiguity Classification
Systems}\label{ambiguity-classification-systems}

Our ambiguity detection framework extends beyond simple polysemy to
include context-dependent meaning shifts characteristic of scientific
terminology evolution:

\textbf{Multi-Level Ambiguity Classification}:

\begin{enumerate}
\item **Lexical Ambiguity**: Multiple dictionary meanings (e.g., "individual" in biological vs. psychological contexts)
\item **Contextual Ambiguity**: Meaning shifts based on research tradition (e.g., "caste" in classical vs. modern entomology)
\item **Scale Ambiguity**: Meaning variations across biological scales (e.g., "behavior" at individual vs. colony levels)
\item **Temporal Ambiguity**: Historical meaning evolution (e.g., "superorganism" from 1970s to present)
\end{enumerate}

\subsubsection{Cross-Domain Conceptual
Mapping}\label{cross-domain-conceptual-mapping}

We develop conceptual mapping techniques that reveal relationships
between domains:

\begin{equation}\label{eq:cross_domain_mapping}
M_{ij} = \frac{1}{|T_i \cap T_j|} \sum_{t \in T_i \cap T_j} s(t, D_i, D_j)
\end{equation}

where \(M_{ij}\) is the mapping strength between domains \(D_i\) and
\(D_j\), and \(s(t, D_i, D_j)\) measures semantic similarity of term
\(t\) across domains.

\subsection{Framing Analysis Methods}\label{framing-analysis-methods}

\subsubsection{Anthropomorphic Framing
Detection}\label{anthropomorphic-framing-detection}

Anthropomorphic framing detection incorporates linguistic and conceptual
indicators:

\textbf{Linguistic Indicators}:

\begin{itemize}
\tightlist
\item
  Pronominalization (use of ``it'' vs.~``she/he'' for colonies)
\item
  Agency attribution (active vs.~passive voice patterns)
\item
  Intentionality markers (words implying purpose or planning)
\end{itemize}

\textbf{Conceptual Indicators}:

\begin{itemize}
\tightlist
\item
  Social structure projections (human hierarchies onto insect societies)
\item
  Emotional attribution (anthropomorphic emotional terms)
\item
  Cultural bias patterns (Western social norms in biological
  descriptions)
\end{itemize}

\subsubsection{Hierarchical Framing
Analysis}\label{hierarchical-framing-analysis}

Analysis of hierarchical framing reveals nested levels of social
structure imposition:

\textbf{Macro-Level Hierarchies}: Colony-level social organization
(queen → workers → males)

\textbf{Micro-Level Hierarchies}: Individual-level interactions
(dominant → subordinate nestmates)

\textbf{Inter-Colony Hierarchies}: Population-level relationships
(territorial dominance, resource competition)

\subsection{Network Analysis}\label{network-analysis}

\subsubsection{Temporal Network
Evolution}\label{temporal-network-evolution}

Analysis of how terminology networks evolve over time reveals conceptual
shifts:

\begin{equation}\label{eq:temporal_network_evolution}
\Delta G(t) = G(t+1) - G(t) = \sum_{e \in E} \delta_e(t) + \sum_{v \in V} \delta_v(t)
\end{equation}

where \(\delta_e(t)\) and \(\delta_v(t)\) represent edge and vertex
changes over time periods.

\textbf{Key Evolutionary Patterns}:

\begin{itemize}
\tightlist
\item
  \textbf{Network Growth}: Addition of new terms and relationships
\item
  \textbf{Structural Rearrangements}: Changes in network topology
\item
  \textbf{Conceptual Consolidation}: Strengthening of established
  relationships
\item
  \textbf{Paradigm Shifts}: Major restructuring events
\end{itemize}

\subsubsection{Multi-Scale Network
Analysis}\label{multi-scale-network-analysis}

Network analysis at multiple scales reveals hierarchical organization:

\textbf{Local Scale}: Individual term relationships within domains
\textbf{Domain Scale}: Inter-term relationships within domains
\textbf{Cross-Domain Scale}: Relationships between domains \textbf{Field
Scale}: Relationships across the entire entomological terminology
network

\newpage

\section{Supplemental Analysis: Case Studies and
Validation}\label{sec:supplemental_case_studies}

\subsection{Validation Frameworks}\label{validation-frameworks}

\subsubsection{Inter-Subjectivity
Validation}\label{inter-subjectivity-validation}

Validation incorporates multiple perspectives:

\textbf{Expert Validation}: Entomological domain experts review
classifications \textbf{Peer Validation}: Interdisciplinary researchers
assess cross-domain mappings \textbf{Historical Validation}: Analysis of
terminology evolution against known conceptual shifts
\textbf{Cross-Cultural Validation}: Comparison with non-English
entomological literature

\subsubsection{Robustness Testing}\label{robustness-testing}

Robustness analysis ensures result stability:

\textbf{Subsampling Stability}: Performance across different corpus
subsets \textbf{Parameter Sensitivity}: Robustness to algorithmic
parameter variations \textbf{Annotation Consistency}: Agreement across
multiple human annotators \textbf{Temporal Stability}: Consistency
across publication periods

\subsection{Case Study Analysis}\label{case-study-analysis}

\subsubsection{Caste Terminology Evolution:
1850-2024}\label{caste-terminology-evolution-1850-2024}

Ultra-longitudinal analysis reveals century-scale conceptual evolution:

\textbf{Pre-Darwinian Period (1850-1859)}: Essentialist caste categories
based on morphological differences

\textbf{Darwinian Synthesis (1860-1899)}: Evolutionary explanations for
caste differences

\textbf{Genetic Revolution (1900-1949)}: Chromosomal mechanisms
underlying caste determination

\textbf{Molecular Biology Era (1950-1999)}: Gene expression and hormonal
control of caste differentiation

\textbf{Genomic Era (2000-2024)}: Epigenetic and transcriptomic
regulation of caste phenotypes \cite{chandra2021epigenetics},
accompanied by growing recognition that rigid caste categories fail to
capture the labile, environmentally responsive nature of social insect
development \cite{boomsma2018superorganismality}.
\citet{warner2024caste} demonstrate that caste differentiation becomes
increasingly \emph{canalized} from early development through cascading
gene-expression changes modulated by juvenile hormone signaling, while
gene expression in \emph{Lasius niger} is more strongly influenced by
age than by caste---further undermining the fixedness implied by
``caste'' terminology.

\paragraph{BHL Historical-Layer Grounding
(1850-1970)}\label{sec:bhl_grounding}

These period claims are anchored in measurable term usage by the
Biodiversity Heritage Library historical full-text layer: 2460
ant/myrmecology documents harvested from the BHL mirror collection on
the Internet Archive (BHL API v3 requires a key; see
\texttt{data/bhl/README.md} for the provenance and the verbatim query).
The era-stratified artifact \texttt{data/bhl/era\_term\_usage.json},
produced by \texttt{src/pipeline/bhl\_analysis.py}, reports ``caste''
rising monotonically from 0.0836 per 10k tokens in 1850-1899 (1013
documents) to 0.1130 in 1900-1949 (1247 documents) and 0.2530 in
1950-1970 (200 documents) --- an intensification consistent with the
Genetic Revolution and Molecular Biology era framings above. The labor
vocabulary shows the same arc: ``worker'' rises from 0.6583 to 0.7938 to
1.4608 per 10k tokens across the eras, while ``queen'' remains the
anchor term throughout (0.8993 → 1.0799 → 1.1384 per 10k). The
integrative-era vocabulary is absent from the entire stratum:
``superorganism'' registers 0.0017, 0.0008, and 0.0099 per 10k tokens
across the three eras --- zero throughout --- marking it as a post-1970
import into general monographic usage rather than a term the historical
corpus itself carries. All numbers resolve at render time from the
BHL\_* manuscript tokens.

\subsubsection{Superorganism Concept
Evolution}\label{superorganism-concept-evolution}

Table \ref{tab:superorganism_concept_evolution} traces the superorganism
concept across seven decades of research:

The BHL historical layer brackets the start of this table from below:
across 2460 BHL-mirror documents, ``superorganism'' occurs 0.0017 per
10k tokens in 1850-1899, 0.0008 in 1900-1949, and 0.0099 in 1950-1970
--- zero throughout (\texttt{data/bhl/era\_term\_usage.json}, produced
by \texttt{src/pipeline/bhl\_analysis.py}), confirming that the
concept's theoretical vocabulary lived outside the monographic stratum
and entered routine usage only in the post-1970 decades traced here.

\begin{table}[h]
\centering
\begin{tabular}{|l|c|c|c|c|}
\hline
\textbf{Era} & \textbf{Dominant Metaphor} & \textbf{Key Evidence} & \textbf{Critiques} & \textbf{Legacy} \\
\hline
1960s & Organismic & Division of labor analogies & Ignores individual variation & Established field \\
1970s & Cybernetic & Communication networks & Mechanistic reductionism & Systems thinking \\
1980s & Genetic & Kin selection theory & Haplodiploidy focus & Evolutionary framework \\
1990s & Neuroendocrine & Pheromonal control & Colony complexity & Regulatory mechanisms \\
2000s & Epigenetic & DNA methylation & Environmental effects & Developmental plasticity \\
2010s & Microbiome & Symbiont communities & Host-symbiont dynamics & Extended organism concept \\
2020s & Canalization & Cascading gene expression & Lability of ``caste'' & Terminological reform \\
\hline
\end{tabular}
\caption{Evolution of superorganism concept across research eras}
\label{tab:superorganism_concept_evolution}
\end{table}

\subsection{Methodological
Reflections}\label{methodological-reflections}

\subsubsection{Mixed-Methodology
Integration}\label{mixed-methodology-integration}

Our approach integrates qualitative and quantitative methods:

\textbf{Qualitative Contributions}:

\begin{itemize}
\tightlist
\item
  Theoretical framework development
\item
  Conceptual category identification
\item
  Historical context analysis
\item
  Cross-domain relationship mapping
\end{itemize}

\textbf{Quantitative Contributions}:

\begin{itemize}
\tightlist
\item
  Statistical pattern identification
\item
  Network structure analysis
\item
  Temporal trend quantification
\item
  Validation metric development
\end{itemize}

For a discussion of methodological limitations and scope considerations,
see Section \ref{sec:discussion}. Future research directions, including
semantic analysis (transformer-based embeddings, multilingual
extensions) and practical applications (terminology standards, peer
review tools), are discussed in Section \ref{sec:conclusion}.

\newpage

\end{document}
